\documentclass[
  doc,
  floatsintext,
  longtable,
  a4paper,
  nolmodern,
  notxfonts,
  notimes,
  12pt,
  colorlinks=true,linkcolor=blue,citecolor=blue,urlcolor=blue]{apa7}

\usepackage{amsmath}
\usepackage{amssymb}

\geometry{inner=1in, outer=1in}
\fancyhfoffset[LE,RO]{0cm}


\usepackage[bidi=default]{babel}
\babelprovide[main,import]{spanish}


% get rid of language-specific shorthands (see #6817):
\let\LanguageShortHands\languageshorthands
\def\languageshorthands#1{}

\RequirePackage{longtable}
\RequirePackage{threeparttablex}

\makeatletter
\renewcommand{\paragraph}{\@startsection{paragraph}{4}{\parindent}%
	{0\baselineskip \@plus 0.2ex \@minus 0.2ex}%
	{-.5em}%
	{\normalfont\normalsize\bfseries\typesectitle}}

\renewcommand{\subparagraph}[1]{\@startsection{subparagraph}{5}{0.5em}%
	{0\baselineskip \@plus 0.2ex \@minus 0.2ex}%
	{-\z@\relax}%
	{\normalfont\normalsize\bfseries\itshape\hspace{\parindent}{#1}\textit{\addperi}}{\relax}}
\makeatother




\usepackage{longtable, booktabs, multirow, multicol, colortbl, hhline, caption, array, float, xpatch}
\usepackage{subcaption}


\renewcommand\thesubfigure{\Alph{subfigure}}
\setcounter{topnumber}{2}
\setcounter{bottomnumber}{2}
\setcounter{totalnumber}{4}
\renewcommand{\topfraction}{0.85}
\renewcommand{\bottomfraction}{0.85}
\renewcommand{\textfraction}{0.15}
\renewcommand{\floatpagefraction}{0.7}

\usepackage{tcolorbox}
\tcbuselibrary{listings,theorems, breakable, skins}
\usepackage{fontawesome5}

\definecolor{quarto-callout-color}{HTML}{909090}
\definecolor{quarto-callout-note-color}{HTML}{0758E5}
\definecolor{quarto-callout-important-color}{HTML}{CC1914}
\definecolor{quarto-callout-warning-color}{HTML}{EB9113}
\definecolor{quarto-callout-tip-color}{HTML}{00A047}
\definecolor{quarto-callout-caution-color}{HTML}{FC5300}
\definecolor{quarto-callout-color-frame}{HTML}{ACACAC}
\definecolor{quarto-callout-note-color-frame}{HTML}{4582EC}
\definecolor{quarto-callout-important-color-frame}{HTML}{D9534F}
\definecolor{quarto-callout-warning-color-frame}{HTML}{F0AD4E}
\definecolor{quarto-callout-tip-color-frame}{HTML}{02B875}
\definecolor{quarto-callout-caution-color-frame}{HTML}{FD7E14}

%\newlength\Oldarrayrulewidth
%\newlength\Oldtabcolsep


\usepackage{hyperref}



\usepackage{color}
\usepackage{fancyvrb}
\newcommand{\VerbBar}{|}
\newcommand{\VERB}{\Verb[commandchars=\\\{\}]}
\DefineVerbatimEnvironment{Highlighting}{Verbatim}{commandchars=\\\{\}}
% Add ',fontsize=\small' for more characters per line
\usepackage{framed}
\definecolor{shadecolor}{RGB}{241,243,245}
\newenvironment{Shaded}{\begin{snugshade}}{\end{snugshade}}
\newcommand{\AlertTok}[1]{\textcolor[rgb]{0.68,0.00,0.00}{#1}}
\newcommand{\AnnotationTok}[1]{\textcolor[rgb]{0.37,0.37,0.37}{#1}}
\newcommand{\AttributeTok}[1]{\textcolor[rgb]{0.40,0.45,0.13}{#1}}
\newcommand{\BaseNTok}[1]{\textcolor[rgb]{0.68,0.00,0.00}{#1}}
\newcommand{\BuiltInTok}[1]{\textcolor[rgb]{0.00,0.23,0.31}{#1}}
\newcommand{\CharTok}[1]{\textcolor[rgb]{0.13,0.47,0.30}{#1}}
\newcommand{\CommentTok}[1]{\textcolor[rgb]{0.37,0.37,0.37}{#1}}
\newcommand{\CommentVarTok}[1]{\textcolor[rgb]{0.37,0.37,0.37}{\textit{#1}}}
\newcommand{\ConstantTok}[1]{\textcolor[rgb]{0.56,0.35,0.01}{#1}}
\newcommand{\ControlFlowTok}[1]{\textcolor[rgb]{0.00,0.23,0.31}{\textbf{#1}}}
\newcommand{\DataTypeTok}[1]{\textcolor[rgb]{0.68,0.00,0.00}{#1}}
\newcommand{\DecValTok}[1]{\textcolor[rgb]{0.68,0.00,0.00}{#1}}
\newcommand{\DocumentationTok}[1]{\textcolor[rgb]{0.37,0.37,0.37}{\textit{#1}}}
\newcommand{\ErrorTok}[1]{\textcolor[rgb]{0.68,0.00,0.00}{#1}}
\newcommand{\ExtensionTok}[1]{\textcolor[rgb]{0.00,0.23,0.31}{#1}}
\newcommand{\FloatTok}[1]{\textcolor[rgb]{0.68,0.00,0.00}{#1}}
\newcommand{\FunctionTok}[1]{\textcolor[rgb]{0.28,0.35,0.67}{#1}}
\newcommand{\ImportTok}[1]{\textcolor[rgb]{0.00,0.46,0.62}{#1}}
\newcommand{\InformationTok}[1]{\textcolor[rgb]{0.37,0.37,0.37}{#1}}
\newcommand{\KeywordTok}[1]{\textcolor[rgb]{0.00,0.23,0.31}{\textbf{#1}}}
\newcommand{\NormalTok}[1]{\textcolor[rgb]{0.00,0.23,0.31}{#1}}
\newcommand{\OperatorTok}[1]{\textcolor[rgb]{0.37,0.37,0.37}{#1}}
\newcommand{\OtherTok}[1]{\textcolor[rgb]{0.00,0.23,0.31}{#1}}
\newcommand{\PreprocessorTok}[1]{\textcolor[rgb]{0.68,0.00,0.00}{#1}}
\newcommand{\RegionMarkerTok}[1]{\textcolor[rgb]{0.00,0.23,0.31}{#1}}
\newcommand{\SpecialCharTok}[1]{\textcolor[rgb]{0.37,0.37,0.37}{#1}}
\newcommand{\SpecialStringTok}[1]{\textcolor[rgb]{0.13,0.47,0.30}{#1}}
\newcommand{\StringTok}[1]{\textcolor[rgb]{0.13,0.47,0.30}{#1}}
\newcommand{\VariableTok}[1]{\textcolor[rgb]{0.07,0.07,0.07}{#1}}
\newcommand{\VerbatimStringTok}[1]{\textcolor[rgb]{0.13,0.47,0.30}{#1}}
\newcommand{\WarningTok}[1]{\textcolor[rgb]{0.37,0.37,0.37}{\textit{#1}}}

\providecommand{\tightlist}{%
  \setlength{\itemsep}{0pt}\setlength{\parskip}{0pt}}
\usepackage{longtable,booktabs,array}
\usepackage{calc} % for calculating minipage widths
% Correct order of tables after \paragraph or \subparagraph
\usepackage{etoolbox}
\makeatletter
\patchcmd\longtable{\par}{\if@noskipsec\mbox{}\fi\par}{}{}
\makeatother
% Allow footnotes in longtable head/foot
\IfFileExists{footnotehyper.sty}{\usepackage{footnotehyper}}{\usepackage{footnote}}
\makesavenoteenv{longtable}

\usepackage{graphicx}
\makeatletter
\newsavebox\pandoc@box
\newcommand*\pandocbounded[1]{% scales image to fit in text height/width
  \sbox\pandoc@box{#1}%
  \Gscale@div\@tempa{\textheight}{\dimexpr\ht\pandoc@box+\dp\pandoc@box\relax}%
  \Gscale@div\@tempb{\linewidth}{\wd\pandoc@box}%
  \ifdim\@tempb\p@<\@tempa\p@\let\@tempa\@tempb\fi% select the smaller of both
  \ifdim\@tempa\p@<\p@\scalebox{\@tempa}{\usebox\pandoc@box}%
  \else\usebox{\pandoc@box}%
  \fi%
}
% Set default figure placement to htbp
\def\fps@figure{htbp}
\makeatother







\usepackage{newtx}

\defaultfontfeatures{Scale=MatchLowercase}
\defaultfontfeatures[\rmfamily]{Ligatures=TeX,Scale=1}





\title{Modelos de regresión en Machine Learning}


\shorttitle{ML PARA REGRESIÓN}


\usepackage{etoolbox}



\ccoppy{\textcopyright~2021}



\author{Edison Achalma}



\affiliation{
{Escuela Profesional de Economía, Universidad Nacional de San Cristóbal
de Huamanga}}




\leftheader{Achalma}

\date{2021-12-13}


\abstract{Este abstract será actualizado una vez que se complete el
contenido final del artículo. }

\keywords{keyword1, keyword2}

\authornote{\par{\addORCIDlink{Edison Achalma}{0000-0001-6996-3364}} 
\par{ }
\par{   El autor no tiene conflictos de interés que revelar.    Los
roles de autor se clasificaron utilizando la taxonomía de roles de
colaborador (CRediT; https://credit.niso.org/) de la siguiente
manera:  Edison Achalma:   conceptualización, metodología, análisis
formal, investigación, recursos, curación de
datos, redacción, visualización, supervisión, administración del
proyecto}
\par{La correspondencia relativa a este artículo debe dirigirse a Edison
Achalma, Escuela Profesional de Economía, Universidad Nacional de San
Cristóbal de Huamanga, Portal Independencia N°
57, Ayacucho, AYA 5001, Perú, Email: \href{mailto:elmer.achalma.09@unsch.edu.pe}{elmer.achalma.09@unsch.edu.pe}}
}

\makeatletter
\let\endoldlt\endlongtable
\def\endlongtable{
\hline
\endoldlt
}
\makeatother

\urlstyle{same}



\makeatletter
\@ifpackageloaded{caption}{}{\usepackage{caption}}
\AtBeginDocument{%
\ifdefined\contentsname
  \renewcommand*\contentsname{Tabla de contenidos}
\else
  \newcommand\contentsname{Tabla de contenidos}
\fi
\ifdefined\listfigurename
  \renewcommand*\listfigurename{Lista de Figuras}
\else
  \newcommand\listfigurename{Lista de Figuras}
\fi
\ifdefined\listtablename
  \renewcommand*\listtablename{Lista de Tablas}
\else
  \newcommand\listtablename{Lista de Tablas}
\fi
\ifdefined\figurename
  \renewcommand*\figurename{Figura}
\else
  \newcommand\figurename{Figura}
\fi
\ifdefined\tablename
  \renewcommand*\tablename{Tabla}
\else
  \newcommand\tablename{Tabla}
\fi
}
\@ifpackageloaded{float}{}{\usepackage{float}}
\floatstyle{ruled}
\@ifundefined{c@chapter}{\newfloat{codelisting}{h}{lop}}{\newfloat{codelisting}{h}{lop}[chapter]}
\floatname{codelisting}{Listado}
\newcommand*\listoflistings{\listof{codelisting}{Lista de Listados}}
\makeatother
\makeatletter
\makeatother
\makeatletter
\@ifpackageloaded{caption}{}{\usepackage{caption}}
\@ifpackageloaded{subcaption}{}{\usepackage{subcaption}}
\makeatother
\makeatletter
\@ifpackageloaded{fontawesome5}{}{\usepackage{fontawesome5}}
\makeatother

% From https://tex.stackexchange.com/a/645996/211326
%%% apa7 doesn't want to add appendix section titles in the toc
%%% let's make it do it
\makeatletter
\xpatchcmd{\appendix}
  {\par}
  {\addcontentsline{toc}{section}{\@currentlabelname}\par}
  {}{}
\makeatother

%% Disable longtable counter
%% https://tex.stackexchange.com/a/248395/211326

\usepackage{etoolbox}

\makeatletter
\patchcmd{\LT@caption}
  {\bgroup}
  {\bgroup\global\LTpatch@captiontrue}
  {}{}
\patchcmd{\longtable}
  {\par}
  {\par\global\LTpatch@captionfalse}
  {}{}
\apptocmd{\endlongtable}
  {\ifLTpatch@caption\else\addtocounter{table}{-1}\fi}
  {}{}
\newif\ifLTpatch@caption
\makeatother

\begin{document}

\maketitle


\hypertarget{toc}{}
\tableofcontents
\newpage
\section[Introduction]{Modelos de regresión en Machine Learning}

\setcounter{secnumdepth}{3}

\setlength\LTleft{0pt}




En esta décima guía exploraremos los principales algoritmos de machine
learning para problemas de regresión. Estos métodos son fundamentales
para resolver problemas económicos como predicción de precios,
estimación de demanda, forecasting de series de tiempo, valuación de
activos, entre otros.

\section{La regresión}\label{la-regresiuxf3n}

\subsection{Concepto}\label{concepto}

La regresión es un método para predecir valores numéricos continuos (a
diferencia de clasificación que predice categorías discretas).

\begin{Shaded}
\begin{Highlighting}[]
\ImportTok{import}\NormalTok{ pandas }\ImportTok{as}\NormalTok{ pd}
\ImportTok{import}\NormalTok{ numpy }\ImportTok{as}\NormalTok{ np}
\ImportTok{import}\NormalTok{ matplotlib.pyplot }\ImportTok{as}\NormalTok{ plt}
\ImportTok{import}\NormalTok{ seaborn }\ImportTok{as}\NormalTok{ sns}

\CommentTok{\# Scikit{-}learn}
\ImportTok{from}\NormalTok{ sklearn.model\_selection }\ImportTok{import}\NormalTok{ train\_test\_split, cross\_val\_score}
\ImportTok{from}\NormalTok{ sklearn.preprocessing }\ImportTok{import}\NormalTok{ StandardScaler, PolynomialFeatures}
\ImportTok{from}\NormalTok{ sklearn.linear\_model }\ImportTok{import}\NormalTok{ LinearRegression, Ridge, Lasso, ElasticNet}
\ImportTok{from}\NormalTok{ sklearn.neighbors }\ImportTok{import}\NormalTok{ KNeighborsRegressor}
\ImportTok{from}\NormalTok{ sklearn.tree }\ImportTok{import}\NormalTok{ DecisionTreeRegressor}
\ImportTok{from}\NormalTok{ sklearn.ensemble }\ImportTok{import}\NormalTok{ RandomForestRegressor}
\ImportTok{from}\NormalTok{ sklearn.metrics }\ImportTok{import}\NormalTok{ mean\_squared\_error, mean\_absolute\_error, r2\_score}

\CommentTok{\# Configuración}
\NormalTok{plt.style.use(}\StringTok{\textquotesingle{}seaborn{-}v0\_8{-}whitegrid\textquotesingle{}}\NormalTok{)}
\NormalTok{plt.rcParams[}\StringTok{\textquotesingle{}figure.figsize\textquotesingle{}}\NormalTok{] }\OperatorTok{=}\NormalTok{ (}\DecValTok{12}\NormalTok{, }\DecValTok{6}\NormalTok{)}
\NormalTok{np.random.seed(}\DecValTok{42}\NormalTok{)}

\BuiltInTok{print}\NormalTok{(}\StringTok{"Librerías importadas exitosamente"}\NormalTok{)}
\end{Highlighting}
\end{Shaded}

\subsection{Tipos de problemas de
regresión}\label{tipos-de-problemas-de-regresiuxf3n}

\begin{enumerate}
\def\labelenumi{\arabic{enumi}.}
\item
  \textbf{Predicción de precios}

  \begin{itemize}
  \tightlist
  \item
    Precios de viviendas
  \item
    Precios de acciones
  \item
    Precios de commodities
  \item
    Tipos de cambio
  \end{itemize}
\item
  \textbf{Estimación de demanda}*

  \begin{itemize}
  \tightlist
  \item
    Demanda de productos
  \item
    Demanda laboral
  \item
    Demanda de servicios públicos
  \end{itemize}
\item
  \textbf{Forecasting económico} (pronósticos)

  \begin{itemize}
  \tightlist
  \item
    PIB futuro
  \item
    Ventas futuras
  \item
    Inflación
  \item
    Tasas de interés
  \end{itemize}
\item
  \textbf{Valuación}

  \begin{itemize}
  \tightlist
  \item
    Valuación de empresas
  \item
    Valuación de activos
  \item
    Riesgo de crédito (score continuo)
  \end{itemize}
\item
  \textbf{Análisis de impacto y causalidad}

  \begin{itemize}
  \tightlist
  \item
    Efecto de políticas públicas
  \item
    Retorno de inversión en publicidad
  \item
    Elasticidades (precio, ingreso)
  \end{itemize}
\end{enumerate}

\subsection{Generar datos de ejemplo}\label{generar-datos-de-ejemplo}

\begin{Shaded}
\begin{Highlighting}[]
\KeywordTok{def}\NormalTok{ generar\_datos\_regresion(n}\OperatorTok{=}\DecValTok{200}\NormalTok{, ruido}\OperatorTok{=}\FloatTok{1.0}\NormalTok{, seed}\OperatorTok{=}\DecValTok{42}\NormalTok{):}
    \CommentTok{"""}
\CommentTok{    Genera datos sintéticos para regresión}
\CommentTok{    }
\CommentTok{    Modelo: Y = 3 + 2*X1 {-} 1.5*X2 + ε}
\CommentTok{    """}
\NormalTok{    np.random.seed(seed)}
    
    \CommentTok{\# Variables explicativas}
\NormalTok{    X }\OperatorTok{=}\NormalTok{ np.random.randn(n, }\DecValTok{2}\NormalTok{)}
    
    \CommentTok{\# Coeficientes verdaderos}
\NormalTok{    beta\_true }\OperatorTok{=}\NormalTok{ np.array([}\FloatTok{2.0}\NormalTok{, }\OperatorTok{{-}}\FloatTok{1.5}\NormalTok{])}
\NormalTok{    intercepto\_true }\OperatorTok{=} \FloatTok{3.0}
    
    \CommentTok{\# Variable dependiente con ruido}
\NormalTok{    Y }\OperatorTok{=}\NormalTok{ intercepto\_true }\OperatorTok{+}\NormalTok{ np.dot(X, beta\_true) }\OperatorTok{+}\NormalTok{ np.random.randn(n) }\OperatorTok{*}\NormalTok{ ruido}
    
    \ControlFlowTok{return}\NormalTok{ X, Y, beta\_true, intercepto\_true}

\CommentTok{\# Generar datos}
\NormalTok{X, Y, beta\_true, intercepto\_true }\OperatorTok{=}\NormalTok{ generar\_datos\_regresion(n}\OperatorTok{=}\DecValTok{200}\NormalTok{, ruido}\OperatorTok{=}\FloatTok{1.0}\NormalTok{)}

\BuiltInTok{print}\NormalTok{(}\StringTok{"DATOS GENERADOS"}\NormalTok{)}
\BuiltInTok{print}\NormalTok{(}\StringTok{"="} \OperatorTok{*} \DecValTok{70}\NormalTok{)}
\BuiltInTok{print}\NormalTok{(}\SpecialStringTok{f"Número de observaciones: }\SpecialCharTok{\{}\BuiltInTok{len}\NormalTok{(Y)}\SpecialCharTok{\}}\SpecialStringTok{"}\NormalTok{)}
\BuiltInTok{print}\NormalTok{(}\SpecialStringTok{f"Número de características: }\SpecialCharTok{\{}\NormalTok{X}\SpecialCharTok{.}\NormalTok{shape[}\DecValTok{1}\NormalTok{]}\SpecialCharTok{\}}\SpecialStringTok{"}\NormalTok{)}
\BuiltInTok{print}\NormalTok{(}\SpecialStringTok{f"}\CharTok{\textbackslash{}n}\SpecialStringTok{Coeficientes verdaderos:"}\NormalTok{)}
\BuiltInTok{print}\NormalTok{(}\SpecialStringTok{f"  Intercepto: }\SpecialCharTok{\{}\NormalTok{intercepto\_true}\SpecialCharTok{\}}\SpecialStringTok{"}\NormalTok{)}
\BuiltInTok{print}\NormalTok{(}\SpecialStringTok{f"  β₁: }\SpecialCharTok{\{}\NormalTok{beta\_true[}\DecValTok{0}\NormalTok{]}\SpecialCharTok{\}}\SpecialStringTok{"}\NormalTok{)}
\BuiltInTok{print}\NormalTok{(}\SpecialStringTok{f"  β₂: }\SpecialCharTok{\{}\NormalTok{beta\_true[}\DecValTok{1}\NormalTok{]}\SpecialCharTok{\}}\SpecialStringTok{"}\NormalTok{)}
\BuiltInTok{print}\NormalTok{(}\SpecialStringTok{f"}\CharTok{\textbackslash{}n}\SpecialStringTok{Modelo verdadero: Y = }\SpecialCharTok{\{}\NormalTok{intercepto\_true}\SpecialCharTok{\}}\SpecialStringTok{ + }\SpecialCharTok{\{}\NormalTok{beta\_true[}\DecValTok{0}\NormalTok{]}\SpecialCharTok{\}}\SpecialStringTok{*X₁ + }\SpecialCharTok{\{}\NormalTok{beta\_true[}\DecValTok{1}\NormalTok{]}\SpecialCharTok{\}}\SpecialStringTok{*X₂ + ε"}\NormalTok{)}

\CommentTok{\# Visualizar datos}
\NormalTok{fig, axes }\OperatorTok{=}\NormalTok{ plt.subplots(}\DecValTok{1}\NormalTok{, }\DecValTok{2}\NormalTok{, figsize}\OperatorTok{=}\NormalTok{(}\DecValTok{14}\NormalTok{, }\DecValTok{5}\NormalTok{))}

\NormalTok{axes[}\DecValTok{0}\NormalTok{].scatter(X[:, }\DecValTok{0}\NormalTok{], Y, alpha}\OperatorTok{=}\FloatTok{0.6}\NormalTok{, s}\OperatorTok{=}\DecValTok{50}\NormalTok{)}
\NormalTok{axes[}\DecValTok{0}\NormalTok{].set\_xlabel(}\StringTok{\textquotesingle{}X₁\textquotesingle{}}\NormalTok{)}
\NormalTok{axes[}\DecValTok{0}\NormalTok{].set\_ylabel(}\StringTok{\textquotesingle{}Y\textquotesingle{}}\NormalTok{)}
\NormalTok{axes[}\DecValTok{0}\NormalTok{].set\_title(}\StringTok{\textquotesingle{}Relación Y vs X₁\textquotesingle{}}\NormalTok{)}
\NormalTok{axes[}\DecValTok{0}\NormalTok{].grid(}\VariableTok{True}\NormalTok{, alpha}\OperatorTok{=}\FloatTok{0.3}\NormalTok{)}

\NormalTok{axes[}\DecValTok{1}\NormalTok{].scatter(X[:, }\DecValTok{1}\NormalTok{], Y, alpha}\OperatorTok{=}\FloatTok{0.6}\NormalTok{, s}\OperatorTok{=}\DecValTok{50}\NormalTok{, color}\OperatorTok{=}\StringTok{\textquotesingle{}orange\textquotesingle{}}\NormalTok{)}
\NormalTok{axes[}\DecValTok{1}\NormalTok{].set\_xlabel(}\StringTok{\textquotesingle{}X₂\textquotesingle{}}\NormalTok{)}
\NormalTok{axes[}\DecValTok{1}\NormalTok{].set\_ylabel(}\StringTok{\textquotesingle{}Y\textquotesingle{}}\NormalTok{)}
\NormalTok{axes[}\DecValTok{1}\NormalTok{].set\_title(}\StringTok{\textquotesingle{}Relación Y vs X₂\textquotesingle{}}\NormalTok{)}
\NormalTok{axes[}\DecValTok{1}\NormalTok{].grid(}\VariableTok{True}\NormalTok{, alpha}\OperatorTok{=}\FloatTok{0.3}\NormalTok{)}

\NormalTok{plt.tight\_layout()}
\NormalTok{plt.savefig(}\StringTok{\textquotesingle{}datos\_regresion.png\textquotesingle{}}\NormalTok{, dpi}\OperatorTok{=}\DecValTok{300}\NormalTok{, bbox\_inches}\OperatorTok{=}\StringTok{\textquotesingle{}tight\textquotesingle{}}\NormalTok{)}
\BuiltInTok{print}\NormalTok{(}\StringTok{"}\CharTok{\textbackslash{}n}\StringTok{Gráfico guardado como \textquotesingle{}datos\_regresion.png\textquotesingle{}"}\NormalTok{)}
\end{Highlighting}
\end{Shaded}

\subsection{Dividir datos}\label{dividir-datos}

\begin{Shaded}
\begin{Highlighting}[]
\CommentTok{\# Dividir en entrenamiento y prueba}
\NormalTok{X\_train, X\_test, y\_train, y\_test }\OperatorTok{=}\NormalTok{ train\_test\_split(}
\NormalTok{    X, Y, test\_size}\OperatorTok{=}\FloatTok{0.3}\NormalTok{, random\_state}\OperatorTok{=}\DecValTok{42}
\NormalTok{)}

\BuiltInTok{print}\NormalTok{(}\SpecialStringTok{f"}\CharTok{\textbackslash{}n}\SpecialStringTok{División de datos:"}\NormalTok{)}
\BuiltInTok{print}\NormalTok{(}\SpecialStringTok{f"  Entrenamiento: }\SpecialCharTok{\{}\BuiltInTok{len}\NormalTok{(y\_train)}\SpecialCharTok{\}}\SpecialStringTok{ observaciones"}\NormalTok{)}
\BuiltInTok{print}\NormalTok{(}\SpecialStringTok{f"  Prueba: }\SpecialCharTok{\{}\BuiltInTok{len}\NormalTok{(y\_test)}\SpecialCharTok{\}}\SpecialStringTok{ observaciones"}\NormalTok{)}
\end{Highlighting}
\end{Shaded}

\section{Mínimos cuadrados ordinarios
(OLS)}\label{muxednimos-cuadrados-ordinarios-ols}

\subsection{Fundamento matemático}\label{fundamento-matemuxe1tico}

El método de Mínimos Cuadrados Ordinarios (Ordinary Least Squares, OLS)
es el método más fundamental y ampliamente utilizado en econometría y
análisis de regresión. Fue desarrollado por Carl Friedrich Gauss a
principios del siglo XIX y constituye la base para muchos métodos
estadísticos modernos.

\subsubsection{Especificación del
modelo}\label{especificaciuxf3n-del-modelo}

El modelo de regresión lineal múltiple se especifica como:

\[
Y_i = \beta_0 + \beta_1 X_{1i} + \beta_2 X_{2i} + \cdots + \beta_p X_{pi} + \varepsilon_i
\]

donde:

\begin{itemize}
\tightlist
\item
  \(Y_i\) es el valor observado de la variable dependiente para la
  observación \(i\)
\item
  \(X_{ji}\) es el valor de la \(j\)-ésima variable independiente para
  la observación \(i\)
\item
  \(\beta_0\) es el intercepto (término constante)
\item
  \(\beta_j\) son los coeficientes de regresión para
  \(j = 1, 2, \ldots, p\)
\item
  \(\varepsilon_i\) es el término de error aleatorio
\item
  \(i = 1, 2, \ldots, n\) indexa las observaciones
\end{itemize}

En notación matricial, el modelo se expresa de forma más compacta como:

\[
\mathbf{Y} = \mathbf{X}\boldsymbol{\beta} + \boldsymbol{\varepsilon}
\]

donde:

\begin{itemize}
\tightlist
\item
  \(\mathbf{Y}\) es un vector \(n \times 1\) de observaciones de la
  variable dependiente
\item
  \(\mathbf{X}\) es una matriz \(n \times (p+1)\) de variables
  independientes (incluye columna de unos para el intercepto)
\item
  \(\boldsymbol{\beta}\) es un vector \((p+1) \times 1\) de coeficientes
\item
  \(\boldsymbol{\varepsilon}\) es un vector \(n \times 1\) de errores
\end{itemize}

\subsubsection{Función objetivo}\label{funciuxf3n-objetivo}

El método OLS busca estimar los coeficientes \(\boldsymbol{\beta}\) que
minimizan la Suma de Residuos al Cuadrado (Residual Sum of Squares,
RSS):

\[
\text{RSS}(\boldsymbol{\beta}) = \sum_{i=1}^{n} \varepsilon_i^2 = \sum_{i=1}^{n} (Y_i - \hat{Y}_i)^2
\]

donde \(\hat{Y}_i = \beta_0 + \beta_1 X_{1i} + \cdots + \beta_p X_{pi}\)
es el valor predicho.

Expandiendo la expresión:

\[
\text{RSS}(\boldsymbol{\beta}) = \sum_{i=1}^{n} \left(Y_i - \beta_0 - \sum_{j=1}^{p} \beta_j X_{ji}\right)^2
\]

En notación matricial:

\[
\text{RSS}(\boldsymbol{\beta}) = (\mathbf{Y} - \mathbf{X}\boldsymbol{\beta})^\top (\mathbf{Y} - \mathbf{X}\boldsymbol{\beta})
\]

\subsubsection{Derivación de la
solución}\label{derivaciuxf3n-de-la-soluciuxf3n}

Para encontrar el estimador que minimiza RSS, tomamos la derivada con
respecto a \(\boldsymbol{\beta}\) e igualamos a cero:

\[
\frac{\partial \text{RSS}}{\partial \boldsymbol{\beta}} = -2\mathbf{X}^\top(\mathbf{Y} - \mathbf{X}\boldsymbol{\beta}) = 0
\]

Resolviendo para \(\boldsymbol{\beta}\):

\[
\mathbf{X}^\top\mathbf{Y} = \mathbf{X}^\top\mathbf{X}\boldsymbol{\beta}
\]

Estas son las \textbf{ecuaciones normales} del problema de mínimos
cuadrados. Asumiendo que \(\mathbf{X}^\top\mathbf{X}\) es invertible (es
decir, las columnas de \(\mathbf{X}\) son linealmente independientes),
la solución es:

\[
\hat{\boldsymbol{\beta}} = (\mathbf{X}^\top\mathbf{X})^{-1}\mathbf{X}^\top\mathbf{Y}
\]

Esta es la \textbf{solución de mínimos cuadrados ordinarios}.

\subsubsection{Propiedades del estimador
OLS}\label{propiedades-del-estimador-ols}

Bajo los supuestos estándar (ver sección siguiente), el estimador OLS
posee las siguientes propiedades deseables:

\begin{enumerate}
\def\labelenumi{\arabic{enumi}.}
\item
  \textbf{Insesgadez}: El estimador es insesgado, es decir: \[
  E[\hat{\boldsymbol{\beta}}] = \boldsymbol{\beta}
  \]
\item
  \textbf{Varianza mínima}: Según el Teorema de Gauss-Markov, OLS es el
  estimador lineal insesgado de varianza mínima (BLUE: Best Linear
  Unbiased Estimator).
\item
  \textbf{Consistencia}: El estimador converge en probabilidad al
  verdadero valor cuando \(n \to \infty\): \[
  \text{plim}_{n \to \infty} \hat{\boldsymbol{\beta}} = \boldsymbol{\beta}
  \]
\item
  \textbf{Normalidad asintótica}: Bajo condiciones regulares: \[
  \sqrt{n}(\hat{\boldsymbol{\beta}} - \boldsymbol{\beta}) \xrightarrow{d} N(0, \sigma^2(\mathbf{X}^\top\mathbf{X})^{-1})
  \]
\end{enumerate}

\subsubsection{Matriz de
varianzas-covarianzas}\label{matriz-de-varianzas-covarianzas}

La matriz de varianzas-covarianzas del estimador OLS está dada por:

\[
\text{Var}(\hat{\boldsymbol{\beta}}) = \sigma^2 (\mathbf{X}^\top\mathbf{X})^{-1}
\]

donde \(\sigma^2 = \text{Var}(\varepsilon_i)\) es la varianza del error.
En la práctica, \(\sigma^2\) se estima mediante:

\[
\hat{\sigma}^2 = \frac{\text{RSS}}{n - p - 1} = \frac{\sum_{i=1}^{n} \hat{\varepsilon}_i^2}{n - p - 1}
\]

Los errores estándar de los coeficientes son las raíces cuadradas de los
elementos diagonales de
\(\hat{\sigma}^2 (\mathbf{X}^\top\mathbf{X})^{-1}\).

\subsubsection{Supuestos del modelo clásico de regresión
lineal}\label{supuestos-del-modelo-cluxe1sico-de-regresiuxf3n-lineal}

Para que el estimador OLS tenga las propiedades óptimas mencionadas, se
requieren los siguientes supuestos:

\textbf{Supuesto 1: Linealidad}

El modelo es lineal en los parámetros. La verdadera relación entre \(Y\)
y las \(X\) puede expresarse como:

\[
E[Y_i | X_{1i}, \ldots, X_{pi}] = \beta_0 + \beta_1 X_{1i} + \cdots + \beta_p X_{pi}
\]

\textbf{Supuesto 2: Exogeneidad estricta}

Los errores tienen media condicional cero:

\[
E[\varepsilon_i | \mathbf{X}] = 0 \quad \text{para todo } i
\]

Esto implica que las variables independientes no están correlacionadas
con el término de error.

\textbf{Supuesto 3: No colinealidad perfecta}

Las columnas de \(\mathbf{X}\) son linealmente independientes, es decir,
la matriz \(\mathbf{X}^\top\mathbf{X}\) es de rango completo y por tanto
invertible. No existe \(\mathbf{a} \neq \mathbf{0}\) tal que
\(\mathbf{X}\mathbf{a} = \mathbf{0}\).

\textbf{Supuesto 4: Homocedasticidad}

Los errores tienen varianza constante:

\[
\text{Var}(\varepsilon_i | \mathbf{X}) = \sigma^2 \quad \text{para todo } i
\]

\textbf{Supuesto 5: No autocorrelación}

Los errores no están correlacionados entre sí:

\[
\text{Cov}(\varepsilon_i, \varepsilon_j | \mathbf{X}) = 0 \quad \text{para } i \neq j
\]

Los supuestos 4 y 5 pueden combinarse como:

\[
E[\boldsymbol{\varepsilon}\boldsymbol{\varepsilon}^\top | \mathbf{X}] = \sigma^2 \mathbf{I}_n
\]

\textbf{Supuesto 6: Normalidad (opcional)}

Para inferencia en muestras finitas, se asume:

\[
\varepsilon_i | \mathbf{X} \sim N(0, \sigma^2)
\]

Este supuesto no es necesario para propiedades asintóticas pero permite
realizar pruebas de hipótesis exactas en muestras pequeñas.

\subsubsection{Ventajas del método
OLS}\label{ventajas-del-muxe9todo-ols}

\begin{itemize}
\item
  \textbf{Simplicidad conceptual e implementación}: El método tiene una
  interpretación geométrica clara (proyección ortogonal) y es
  computacionalmente eficiente.
\item
  \textbf{Interpretabilidad}: Los coeficientes tienen interpretación
  directa como efectos marginales. \(\beta_j\) representa el cambio en
  \(Y\) por un cambio unitario en \(X_j\), manteniendo constantes las
  demás variables.
\item
  \textbf{Solución cerrada}: Existe una fórmula analítica explícita para
  el estimador, lo que facilita el análisis teórico.
\item
  \textbf{Propiedades estadísticas óptimas}: Bajo los supuestos del
  modelo clásico, OLS es BLUE.
\item
  \textbf{Base para extensiones}: OLS es el fundamento de métodos más
  sofisticados como mínimos cuadrados generalizados (GLS), variables
  instrumentales (IV), y modelos de panel.
\item
  \textbf{Inferencia estadística}: Permite construir intervalos de
  confianza y realizar pruebas de hipótesis sobre los coeficientes.
\end{itemize}

\subsubsection{Limitaciones del método
OLS}\label{limitaciones-del-muxe9todo-ols}

\begin{itemize}
\item
  \textbf{Sensibilidad a violaciones de supuestos}:

  \begin{itemize}
  \tightlist
  \item
    La \textbf{heterocedasticidad} (varianza no constante) invalida los
    errores estándar estándar
  \item
    La \textbf{autocorrelación} (común en series de tiempo) sesga los
    errores estándar
  \item
    La \textbf{endogeneidad} (correlación entre \(X\) y \(\varepsilon\))
    produce estimadores sesgados e inconsistentes
  \end{itemize}
\item
  \textbf{Sensibilidad a outliers}: Valores extremos pueden ejercer
  influencia desproporcionada en las estimaciones debido a la
  minimización de errores al cuadrado.
\item
  \textbf{Multicolinealidad}: Cuando las variables independientes están
  altamente correlacionadas:

  \begin{itemize}
  \tightlist
  \item
    Los errores estándar de los coeficientes se inflan
  \item
    Los coeficientes individuales se vuelven inestables
  \item
    La matriz \(\mathbf{X}^\top\mathbf{X}\) se aproxima a singularidad
  \end{itemize}
\item
  \textbf{Restricción a modelos lineales}: OLS solo puede capturar
  relaciones lineales en los parámetros (aunque puede incluir
  transformaciones no lineales de las variables).
\item
  \textbf{No realiza selección de variables}: Incluye todas las
  variables especificadas sin discriminar su relevancia.
\item
  \textbf{Overfitting con muchas variables}: En modelos con \(p\) grande
  relativo a \(n\), puede ajustar demasiado al ruido de los datos de
  entrenamiento, deteriorando la capacidad predictiva fuera de muestra.
\end{itemize}

\subsubsection{Aplicaciones en
economía}\label{aplicaciones-en-economuxeda}

El método OLS es ubicuo en economía y se aplica a:

\begin{itemize}
\tightlist
\item
  \textbf{Estimación de funciones de demanda}: Relación entre cantidad
  demandada, precio, ingreso y otros factores
\item
  \textbf{Funciones de producción}: Relación entre insumos (capital,
  trabajo) y producto
\item
  \textbf{Ecuaciones de salarios}: Retornos a la educación y experiencia
  (ecuación de Mincer)
\item
  \textbf{Modelos de consumo}: Propensión marginal a consumir
\item
  \textbf{Análisis de series macroeconómicas}: PIB, inflación, desempleo
\item
  \textbf{Evaluación de políticas}: Efectos de programas sociales o
  intervenciones
\end{itemize}

\subsubsection{Consideraciones
computacionales}\label{consideraciones-computacionales}

La inversión de la matriz \(\mathbf{X}^\top\mathbf{X}\) es
computacionalmente costosa cuando \(p\) es grande. Métodos numéricos
estables incluyen:

\begin{itemize}
\tightlist
\item
  \textbf{Descomposición QR}: \(\mathbf{X} = \mathbf{Q}\mathbf{R}\),
  luego
  \(\hat{\boldsymbol{\beta}} = \mathbf{R}^{-1}\mathbf{Q}^\top\mathbf{Y}\)
\item
  \textbf{Descomposición SVD}: Descomposición en valores singulares de
  \(\mathbf{X}\)
\item
  \textbf{Eliminación gaussiana}: Para resolver directamente las
  ecuaciones normales
\end{itemize}

En Python, \texttt{sklearn.linear\_model.LinearRegression} utiliza
descomposición SVD por defecto, lo que proporciona estabilidad numérica
superior.

\subsection{Implementación}\label{implementaciuxf3n}

\begin{Shaded}
\begin{Highlighting}[]
\CommentTok{\# Entrenar modelo OLS}
\NormalTok{modelo\_ols }\OperatorTok{=}\NormalTok{ LinearRegression()}
\NormalTok{modelo\_ols.fit(X\_train, y\_train)}

\CommentTok{\# Coeficientes estimados}
\BuiltInTok{print}\NormalTok{(}\StringTok{"}\CharTok{\textbackslash{}n}\StringTok{MODELO OLS ESTIMADO"}\NormalTok{)}
\BuiltInTok{print}\NormalTok{(}\StringTok{"="} \OperatorTok{*} \DecValTok{70}\NormalTok{)}
\BuiltInTok{print}\NormalTok{(}\SpecialStringTok{f"Intercepto estimado: }\SpecialCharTok{\{}\NormalTok{modelo\_ols}\SpecialCharTok{.}\NormalTok{intercept\_}\SpecialCharTok{:.4f\}}\SpecialStringTok{"}\NormalTok{)}
\BuiltInTok{print}\NormalTok{(}\SpecialStringTok{f"  (Verdadero: }\SpecialCharTok{\{}\NormalTok{intercepto\_true}\SpecialCharTok{\}}\SpecialStringTok{)"}\NormalTok{)}
\BuiltInTok{print}\NormalTok{(}\SpecialStringTok{f"}\CharTok{\textbackslash{}n}\SpecialStringTok{Coeficientes estimados:"}\NormalTok{)}
\BuiltInTok{print}\NormalTok{(}\SpecialStringTok{f"  β₁: }\SpecialCharTok{\{}\NormalTok{modelo\_ols}\SpecialCharTok{.}\NormalTok{coef\_[}\DecValTok{0}\NormalTok{]}\SpecialCharTok{:.4f\}}\SpecialStringTok{ (Verdadero: }\SpecialCharTok{\{}\NormalTok{beta\_true[}\DecValTok{0}\NormalTok{]}\SpecialCharTok{\}}\SpecialStringTok{)"}\NormalTok{)}
\BuiltInTok{print}\NormalTok{(}\SpecialStringTok{f"  β₂: }\SpecialCharTok{\{}\NormalTok{modelo\_ols}\SpecialCharTok{.}\NormalTok{coef\_[}\DecValTok{1}\NormalTok{]}\SpecialCharTok{:.4f\}}\SpecialStringTok{ (Verdadero: }\SpecialCharTok{\{}\NormalTok{beta\_true[}\DecValTok{1}\NormalTok{]}\SpecialCharTok{\}}\SpecialStringTok{)"}\NormalTok{)}

\CommentTok{\# Predecir}
\NormalTok{y\_pred\_ols\_train }\OperatorTok{=}\NormalTok{ modelo\_ols.predict(X\_train)}
\NormalTok{y\_pred\_ols\_test }\OperatorTok{=}\NormalTok{ modelo\_ols.predict(X\_test)}

\CommentTok{\# Evaluar}
\NormalTok{r2\_train }\OperatorTok{=}\NormalTok{ r2\_score(y\_train, y\_pred\_ols\_train)}
\NormalTok{r2\_test }\OperatorTok{=}\NormalTok{ r2\_score(y\_test, y\_pred\_ols\_test)}
\NormalTok{mse\_train }\OperatorTok{=}\NormalTok{ mean\_squared\_error(y\_train, y\_pred\_ols\_train)}
\NormalTok{mse\_test }\OperatorTok{=}\NormalTok{ mean\_squared\_error(y\_test, y\_pred\_ols\_test)}

\BuiltInTok{print}\NormalTok{(}\SpecialStringTok{f"}\CharTok{\textbackslash{}n\textbackslash{}n}\SpecialStringTok{PERFORMANCE DEL MODELO"}\NormalTok{)}
\BuiltInTok{print}\NormalTok{(}\StringTok{"="} \OperatorTok{*} \DecValTok{70}\NormalTok{)}
\BuiltInTok{print}\NormalTok{(}\SpecialStringTok{f"R² Train: }\SpecialCharTok{\{}\NormalTok{r2\_train}\SpecialCharTok{:.4f\}}\SpecialStringTok{"}\NormalTok{)}
\BuiltInTok{print}\NormalTok{(}\SpecialStringTok{f"R² Test:  }\SpecialCharTok{\{}\NormalTok{r2\_test}\SpecialCharTok{:.4f\}}\SpecialStringTok{"}\NormalTok{)}
\BuiltInTok{print}\NormalTok{(}\SpecialStringTok{f"MSE Train: }\SpecialCharTok{\{}\NormalTok{mse\_train}\SpecialCharTok{:.4f\}}\SpecialStringTok{"}\NormalTok{)}
\BuiltInTok{print}\NormalTok{(}\SpecialStringTok{f"MSE Test:  }\SpecialCharTok{\{}\NormalTok{mse\_test}\SpecialCharTok{:.4f\}}\SpecialStringTok{"}\NormalTok{)}
\end{Highlighting}
\end{Shaded}

\subsection{Visualización de
predicciones}\label{visualizaciuxf3n-de-predicciones}

\begin{Shaded}
\begin{Highlighting}[]
\CommentTok{\# Gráfico de valores reales vs predichos}
\NormalTok{fig, axes }\OperatorTok{=}\NormalTok{ plt.subplots(}\DecValTok{1}\NormalTok{, }\DecValTok{2}\NormalTok{, figsize}\OperatorTok{=}\NormalTok{(}\DecValTok{14}\NormalTok{, }\DecValTok{5}\NormalTok{))}

\CommentTok{\# Training set}
\NormalTok{axes[}\DecValTok{0}\NormalTok{].scatter(y\_train, y\_pred\_ols\_train, alpha}\OperatorTok{=}\FloatTok{0.6}\NormalTok{, s}\OperatorTok{=}\DecValTok{50}\NormalTok{)}
\NormalTok{axes[}\DecValTok{0}\NormalTok{].plot([y\_train.}\BuiltInTok{min}\NormalTok{(), y\_train.}\BuiltInTok{max}\NormalTok{()], }
\NormalTok{             [y\_train.}\BuiltInTok{min}\NormalTok{(), y\_train.}\BuiltInTok{max}\NormalTok{()], }
             \StringTok{\textquotesingle{}r{-}{-}\textquotesingle{}}\NormalTok{, linewidth}\OperatorTok{=}\DecValTok{2}\NormalTok{, label}\OperatorTok{=}\StringTok{\textquotesingle{}Predicción perfecta\textquotesingle{}}\NormalTok{)}
\NormalTok{axes[}\DecValTok{0}\NormalTok{].set\_xlabel(}\StringTok{\textquotesingle{}Y real\textquotesingle{}}\NormalTok{)}
\NormalTok{axes[}\DecValTok{0}\NormalTok{].set\_ylabel(}\StringTok{\textquotesingle{}Y predicho\textquotesingle{}}\NormalTok{)}
\NormalTok{axes[}\DecValTok{0}\NormalTok{].set\_title(}\SpecialStringTok{f\textquotesingle{}OLS {-} Entrenamiento (R²=}\SpecialCharTok{\{}\NormalTok{r2\_train}\SpecialCharTok{:.3f\}}\SpecialStringTok{)\textquotesingle{}}\NormalTok{)}
\NormalTok{axes[}\DecValTok{0}\NormalTok{].legend()}
\NormalTok{axes[}\DecValTok{0}\NormalTok{].grid(}\VariableTok{True}\NormalTok{, alpha}\OperatorTok{=}\FloatTok{0.3}\NormalTok{)}

\CommentTok{\# Test set}
\NormalTok{axes[}\DecValTok{1}\NormalTok{].scatter(y\_test, y\_pred\_ols\_test, alpha}\OperatorTok{=}\FloatTok{0.6}\NormalTok{, s}\OperatorTok{=}\DecValTok{50}\NormalTok{, color}\OperatorTok{=}\StringTok{\textquotesingle{}orange\textquotesingle{}}\NormalTok{)}
\NormalTok{axes[}\DecValTok{1}\NormalTok{].plot([y\_test.}\BuiltInTok{min}\NormalTok{(), y\_test.}\BuiltInTok{max}\NormalTok{()], }
\NormalTok{             [y\_test.}\BuiltInTok{min}\NormalTok{(), y\_test.}\BuiltInTok{max}\NormalTok{()], }
             \StringTok{\textquotesingle{}r{-}{-}\textquotesingle{}}\NormalTok{, linewidth}\OperatorTok{=}\DecValTok{2}\NormalTok{, label}\OperatorTok{=}\StringTok{\textquotesingle{}Predicción perfecta\textquotesingle{}}\NormalTok{)}
\NormalTok{axes[}\DecValTok{1}\NormalTok{].set\_xlabel(}\StringTok{\textquotesingle{}Y real\textquotesingle{}}\NormalTok{)}
\NormalTok{axes[}\DecValTok{1}\NormalTok{].set\_ylabel(}\StringTok{\textquotesingle{}Y predicho\textquotesingle{}}\NormalTok{)}
\NormalTok{axes[}\DecValTok{1}\NormalTok{].set\_title(}\SpecialStringTok{f\textquotesingle{}OLS {-} Prueba (R²=}\SpecialCharTok{\{}\NormalTok{r2\_test}\SpecialCharTok{:.3f\}}\SpecialStringTok{)\textquotesingle{}}\NormalTok{)}
\NormalTok{axes[}\DecValTok{1}\NormalTok{].legend()}
\NormalTok{axes[}\DecValTok{1}\NormalTok{].grid(}\VariableTok{True}\NormalTok{, alpha}\OperatorTok{=}\FloatTok{0.3}\NormalTok{)}

\NormalTok{plt.tight\_layout()}
\NormalTok{plt.savefig(}\StringTok{\textquotesingle{}ols\_predictions.png\textquotesingle{}}\NormalTok{, dpi}\OperatorTok{=}\DecValTok{300}\NormalTok{, bbox\_inches}\OperatorTok{=}\StringTok{\textquotesingle{}tight\textquotesingle{}}\NormalTok{)}
\BuiltInTok{print}\NormalTok{(}\StringTok{"}\CharTok{\textbackslash{}n}\StringTok{Gráfico guardado como \textquotesingle{}ols\_predictions.png\textquotesingle{}"}\NormalTok{)}
\end{Highlighting}
\end{Shaded}

\subsection{Análisis de residuos}\label{anuxe1lisis-de-residuos}

\begin{Shaded}
\begin{Highlighting}[]
\CommentTok{\# Calcular residuos}
\NormalTok{residuos\_train }\OperatorTok{=}\NormalTok{ y\_train }\OperatorTok{{-}}\NormalTok{ y\_pred\_ols\_train}
\NormalTok{residuos\_test }\OperatorTok{=}\NormalTok{ y\_test }\OperatorTok{{-}}\NormalTok{ y\_pred\_ols\_test}

\CommentTok{\# Gráficos de diagnóstico}
\NormalTok{fig, axes }\OperatorTok{=}\NormalTok{ plt.subplots(}\DecValTok{2}\NormalTok{, }\DecValTok{2}\NormalTok{, figsize}\OperatorTok{=}\NormalTok{(}\DecValTok{14}\NormalTok{, }\DecValTok{10}\NormalTok{))}

\CommentTok{\# 1. Residuos vs valores predichos}
\NormalTok{axes[}\DecValTok{0}\NormalTok{, }\DecValTok{0}\NormalTok{].scatter(y\_pred\_ols\_train, residuos\_train, alpha}\OperatorTok{=}\FloatTok{0.6}\NormalTok{)}
\NormalTok{axes[}\DecValTok{0}\NormalTok{, }\DecValTok{0}\NormalTok{].axhline(y}\OperatorTok{=}\DecValTok{0}\NormalTok{, color}\OperatorTok{=}\StringTok{\textquotesingle{}r\textquotesingle{}}\NormalTok{, linestyle}\OperatorTok{=}\StringTok{\textquotesingle{}{-}{-}\textquotesingle{}}\NormalTok{, linewidth}\OperatorTok{=}\DecValTok{2}\NormalTok{)}
\NormalTok{axes[}\DecValTok{0}\NormalTok{, }\DecValTok{0}\NormalTok{].set\_xlabel(}\StringTok{\textquotesingle{}Valores predichos\textquotesingle{}}\NormalTok{)}
\NormalTok{axes[}\DecValTok{0}\NormalTok{, }\DecValTok{0}\NormalTok{].set\_ylabel(}\StringTok{\textquotesingle{}Residuos\textquotesingle{}}\NormalTok{)}
\NormalTok{axes[}\DecValTok{0}\NormalTok{, }\DecValTok{0}\NormalTok{].set\_title(}\StringTok{\textquotesingle{}Residuos vs Predicciones\textquotesingle{}}\NormalTok{)}
\NormalTok{axes[}\DecValTok{0}\NormalTok{, }\DecValTok{0}\NormalTok{].grid(}\VariableTok{True}\NormalTok{, alpha}\OperatorTok{=}\FloatTok{0.3}\NormalTok{)}

\CommentTok{\# 2. Histograma de residuos}
\NormalTok{axes[}\DecValTok{0}\NormalTok{, }\DecValTok{1}\NormalTok{].hist(residuos\_train, bins}\OperatorTok{=}\DecValTok{30}\NormalTok{, edgecolor}\OperatorTok{=}\StringTok{\textquotesingle{}black\textquotesingle{}}\NormalTok{, alpha}\OperatorTok{=}\FloatTok{0.7}\NormalTok{)}
\NormalTok{axes[}\DecValTok{0}\NormalTok{, }\DecValTok{1}\NormalTok{].axvline(x}\OperatorTok{=}\DecValTok{0}\NormalTok{, color}\OperatorTok{=}\StringTok{\textquotesingle{}r\textquotesingle{}}\NormalTok{, linestyle}\OperatorTok{=}\StringTok{\textquotesingle{}{-}{-}\textquotesingle{}}\NormalTok{, linewidth}\OperatorTok{=}\DecValTok{2}\NormalTok{)}
\NormalTok{axes[}\DecValTok{0}\NormalTok{, }\DecValTok{1}\NormalTok{].set\_xlabel(}\StringTok{\textquotesingle{}Residuos\textquotesingle{}}\NormalTok{)}
\NormalTok{axes[}\DecValTok{0}\NormalTok{, }\DecValTok{1}\NormalTok{].set\_ylabel(}\StringTok{\textquotesingle{}Frecuencia\textquotesingle{}}\NormalTok{)}
\NormalTok{axes[}\DecValTok{0}\NormalTok{, }\DecValTok{1}\NormalTok{].set\_title(}\StringTok{\textquotesingle{}Distribución de Residuos\textquotesingle{}}\NormalTok{)}
\NormalTok{axes[}\DecValTok{0}\NormalTok{, }\DecValTok{1}\NormalTok{].grid(}\VariableTok{True}\NormalTok{, alpha}\OperatorTok{=}\FloatTok{0.3}\NormalTok{)}

\CommentTok{\# 3. Q{-}Q plot}
\ImportTok{from}\NormalTok{ scipy }\ImportTok{import}\NormalTok{ stats}
\NormalTok{stats.probplot(residuos\_train, dist}\OperatorTok{=}\StringTok{"norm"}\NormalTok{, plot}\OperatorTok{=}\NormalTok{axes[}\DecValTok{1}\NormalTok{, }\DecValTok{0}\NormalTok{])}
\NormalTok{axes[}\DecValTok{1}\NormalTok{, }\DecValTok{0}\NormalTok{].set\_title(}\StringTok{\textquotesingle{}Q{-}Q Plot\textquotesingle{}}\NormalTok{)}
\NormalTok{axes[}\DecValTok{1}\NormalTok{, }\DecValTok{0}\NormalTok{].grid(}\VariableTok{True}\NormalTok{, alpha}\OperatorTok{=}\FloatTok{0.3}\NormalTok{)}

\CommentTok{\# 4. Residuos vs orden}
\NormalTok{axes[}\DecValTok{1}\NormalTok{, }\DecValTok{1}\NormalTok{].scatter(}\BuiltInTok{range}\NormalTok{(}\BuiltInTok{len}\NormalTok{(residuos\_train)), residuos\_train, alpha}\OperatorTok{=}\FloatTok{0.6}\NormalTok{)}
\NormalTok{axes[}\DecValTok{1}\NormalTok{, }\DecValTok{1}\NormalTok{].axhline(y}\OperatorTok{=}\DecValTok{0}\NormalTok{, color}\OperatorTok{=}\StringTok{\textquotesingle{}r\textquotesingle{}}\NormalTok{, linestyle}\OperatorTok{=}\StringTok{\textquotesingle{}{-}{-}\textquotesingle{}}\NormalTok{, linewidth}\OperatorTok{=}\DecValTok{2}\NormalTok{)}
\NormalTok{axes[}\DecValTok{1}\NormalTok{, }\DecValTok{1}\NormalTok{].set\_xlabel(}\StringTok{\textquotesingle{}Índice de observación\textquotesingle{}}\NormalTok{)}
\NormalTok{axes[}\DecValTok{1}\NormalTok{, }\DecValTok{1}\NormalTok{].set\_ylabel(}\StringTok{\textquotesingle{}Residuos\textquotesingle{}}\NormalTok{)}
\NormalTok{axes[}\DecValTok{1}\NormalTok{, }\DecValTok{1}\NormalTok{].set\_title(}\StringTok{\textquotesingle{}Residuos vs Orden\textquotesingle{}}\NormalTok{)}
\NormalTok{axes[}\DecValTok{1}\NormalTok{, }\DecValTok{1}\NormalTok{].grid(}\VariableTok{True}\NormalTok{, alpha}\OperatorTok{=}\FloatTok{0.3}\NormalTok{)}

\NormalTok{plt.tight\_layout()}
\NormalTok{plt.savefig(}\StringTok{\textquotesingle{}ols\_diagnostics.png\textquotesingle{}}\NormalTok{, dpi}\OperatorTok{=}\DecValTok{300}\NormalTok{, bbox\_inches}\OperatorTok{=}\StringTok{\textquotesingle{}tight\textquotesingle{}}\NormalTok{)}
\BuiltInTok{print}\NormalTok{(}\StringTok{"Gráficos de diagnóstico guardados como \textquotesingle{}ols\_diagnostics.png\textquotesingle{}"}\NormalTok{)}
\end{Highlighting}
\end{Shaded}

\section{Regresión Ridge}\label{regresiuxf3n-ridge}

\subsection{Fundamento matemático}\label{fundamento-matemuxe1tico-1}

\begin{Shaded}
\begin{Highlighting}[]
\BuiltInTok{print}\NormalTok{(}\StringTok{"""}
\StringTok{REGRESIÓN RIDGE (L2 REGULARIZATION)}

\StringTok{Objetivo:}
\StringTok{    Minimizar RSS + penalización por magnitud de coeficientes}
\StringTok{    }
\StringTok{    min \{ RSS(β) + α·Σβⱼ² \}}
\StringTok{    }
\StringTok{    = min \{ Σ(yᵢ {-} ŷᵢ)² + α·Σβⱼ² \}}

\StringTok{Solución:}
\StringTok{    β̂ᴿⁱᵈᵍᵉ = (X\textquotesingle{}X + αI)⁻¹X\textquotesingle{}Y}

\StringTok{Donde:}
\StringTok{{-} α: Parámetro de regularización (α ≥ 0)}
\StringTok{  * α = 0: Equivalente a OLS}
\StringTok{  * α → ∞: Coeficientes → 0}
\StringTok{{-} ||β||₂²: Norma L2 (suma de cuadrados de coeficientes)}

\StringTok{Efecto de Ridge:}
\StringTok{✓ Reduce magnitud de coeficientes (shrinkage)}
\StringTok{✓ Mantiene todas las variables (no hace selección)}
\StringTok{✓ Ayuda con multicolinealidad}
\StringTok{✓ Reduce varianza del modelo}
\StringTok{✓ Previene overfitting}

\StringTok{Cuándo usar Ridge:}
\StringTok{{-} Cuando hay muchas variables correlacionadas}
\StringTok{{-} Cuando todas las variables son potencialmente relevantes}
\StringTok{{-} Cuando se prefiere estabilidad sobre interpretabilidad}

\StringTok{Aplicaciones económicas:}
\StringTok{{-} Modelos con muchas variables macroeconómicas correlacionadas}
\StringTok{{-} Predicción de series de tiempo con múltiples rezagos}
\StringTok{{-} Modelos de pricing con características correlacionadas}
\StringTok{"""}\NormalTok{)}
\end{Highlighting}
\end{Shaded}

\subsection{Implementación}\label{implementaciuxf3n-1}

\begin{Shaded}
\begin{Highlighting}[]
\CommentTok{\# Probar diferentes valores de alpha}
\NormalTok{alphas }\OperatorTok{=}\NormalTok{ [}\FloatTok{0.001}\NormalTok{, }\FloatTok{0.01}\NormalTok{, }\FloatTok{0.1}\NormalTok{, }\DecValTok{1}\NormalTok{, }\DecValTok{10}\NormalTok{, }\DecValTok{100}\NormalTok{, }\DecValTok{1000}\NormalTok{]}

\BuiltInTok{print}\NormalTok{(}\StringTok{"}\CharTok{\textbackslash{}n}\StringTok{REGRESIÓN RIDGE {-} EFECTO DE α"}\NormalTok{)}
\BuiltInTok{print}\NormalTok{(}\StringTok{"="} \OperatorTok{*} \DecValTok{70}\NormalTok{)}
\BuiltInTok{print}\NormalTok{(}\SpecialStringTok{f"}\SpecialCharTok{\{}\StringTok{\textquotesingle{}α\textquotesingle{}}\SpecialCharTok{:\textless{}10\}}\SpecialStringTok{ }\SpecialCharTok{\{}\StringTok{\textquotesingle{}β₁\textquotesingle{}}\SpecialCharTok{:\textless{}12\}}\SpecialStringTok{ }\SpecialCharTok{\{}\StringTok{\textquotesingle{}β₂\textquotesingle{}}\SpecialCharTok{:\textless{}12\}}\SpecialStringTok{ }\SpecialCharTok{\{}\StringTok{\textquotesingle{}R² Train\textquotesingle{}}\SpecialCharTok{:\textless{}12\}}\SpecialStringTok{ }\SpecialCharTok{\{}\StringTok{\textquotesingle{}R² Test\textquotesingle{}}\SpecialCharTok{:\textless{}12\}}\SpecialStringTok{"}\NormalTok{)}
\BuiltInTok{print}\NormalTok{(}\StringTok{"{-}"} \OperatorTok{*} \DecValTok{70}\NormalTok{)}

\NormalTok{resultados\_ridge }\OperatorTok{=}\NormalTok{ []}

\ControlFlowTok{for}\NormalTok{ alpha }\KeywordTok{in}\NormalTok{ alphas:}
    \CommentTok{\# Entrenar modelo}
\NormalTok{    modelo\_ridge }\OperatorTok{=}\NormalTok{ Ridge(alpha}\OperatorTok{=}\NormalTok{alpha)}
\NormalTok{    modelo\_ridge.fit(X\_train, y\_train)}
    
    \CommentTok{\# Evaluar}
\NormalTok{    r2\_train }\OperatorTok{=}\NormalTok{ modelo\_ridge.score(X\_train, y\_train)}
\NormalTok{    r2\_test }\OperatorTok{=}\NormalTok{ modelo\_ridge.score(X\_test, y\_test)}
    
\NormalTok{    resultados\_ridge.append(\{}
        \StringTok{\textquotesingle{}alpha\textquotesingle{}}\NormalTok{: alpha,}
        \StringTok{\textquotesingle{}beta1\textquotesingle{}}\NormalTok{: modelo\_ridge.coef\_[}\DecValTok{0}\NormalTok{],}
        \StringTok{\textquotesingle{}beta2\textquotesingle{}}\NormalTok{: modelo\_ridge.coef\_[}\DecValTok{1}\NormalTok{],}
        \StringTok{\textquotesingle{}r2\_train\textquotesingle{}}\NormalTok{: r2\_train,}
        \StringTok{\textquotesingle{}r2\_test\textquotesingle{}}\NormalTok{: r2\_test}
\NormalTok{    \})}
    
    \BuiltInTok{print}\NormalTok{(}\SpecialStringTok{f"}\SpecialCharTok{\{}\NormalTok{alpha}\SpecialCharTok{:\textless{}10.3f\}}\SpecialStringTok{ }\SpecialCharTok{\{}\NormalTok{modelo\_ridge}\SpecialCharTok{.}\NormalTok{coef\_[}\DecValTok{0}\NormalTok{]}\SpecialCharTok{:\textless{}12.4f\}}\SpecialStringTok{ "}
          \SpecialStringTok{f"}\SpecialCharTok{\{}\NormalTok{modelo\_ridge}\SpecialCharTok{.}\NormalTok{coef\_[}\DecValTok{1}\NormalTok{]}\SpecialCharTok{:\textless{}12.4f\}}\SpecialStringTok{ }\SpecialCharTok{\{}\NormalTok{r2\_train}\SpecialCharTok{:\textless{}12.4f\}}\SpecialStringTok{ }\SpecialCharTok{\{}\NormalTok{r2\_test}\SpecialCharTok{:\textless{}12.4f\}}\SpecialStringTok{"}\NormalTok{)}

\NormalTok{df\_ridge }\OperatorTok{=}\NormalTok{ pd.DataFrame(resultados\_ridge)}
\end{Highlighting}
\end{Shaded}

\subsection{Visualización del efecto de
regularización}\label{visualizaciuxf3n-del-efecto-de-regularizaciuxf3n}

\begin{Shaded}
\begin{Highlighting}[]
\CommentTok{\# Gráfico de coeficientes vs alpha}
\NormalTok{fig, axes }\OperatorTok{=}\NormalTok{ plt.subplots(}\DecValTok{1}\NormalTok{, }\DecValTok{2}\NormalTok{, figsize}\OperatorTok{=}\NormalTok{(}\DecValTok{14}\NormalTok{, }\DecValTok{5}\NormalTok{))}

\CommentTok{\# Coeficientes}
\NormalTok{axes[}\DecValTok{0}\NormalTok{].semilogx(df\_ridge[}\StringTok{\textquotesingle{}alpha\textquotesingle{}}\NormalTok{], df\_ridge[}\StringTok{\textquotesingle{}beta1\textquotesingle{}}\NormalTok{], }
                 \StringTok{\textquotesingle{}o{-}\textquotesingle{}}\NormalTok{, linewidth}\OperatorTok{=}\DecValTok{2}\NormalTok{, markersize}\OperatorTok{=}\DecValTok{8}\NormalTok{, label}\OperatorTok{=}\StringTok{\textquotesingle{}β₁\textquotesingle{}}\NormalTok{)}
\NormalTok{axes[}\DecValTok{0}\NormalTok{].semilogx(df\_ridge[}\StringTok{\textquotesingle{}alpha\textquotesingle{}}\NormalTok{], df\_ridge[}\StringTok{\textquotesingle{}beta2\textquotesingle{}}\NormalTok{], }
                 \StringTok{\textquotesingle{}s{-}\textquotesingle{}}\NormalTok{, linewidth}\OperatorTok{=}\DecValTok{2}\NormalTok{, markersize}\OperatorTok{=}\DecValTok{8}\NormalTok{, label}\OperatorTok{=}\StringTok{\textquotesingle{}β₂\textquotesingle{}}\NormalTok{)}
\NormalTok{axes[}\DecValTok{0}\NormalTok{].axhline(y}\OperatorTok{=}\NormalTok{beta\_true[}\DecValTok{0}\NormalTok{], color}\OperatorTok{=}\StringTok{\textquotesingle{}blue\textquotesingle{}}\NormalTok{, linestyle}\OperatorTok{=}\StringTok{\textquotesingle{}{-}{-}\textquotesingle{}}\NormalTok{, }
\NormalTok{                alpha}\OperatorTok{=}\FloatTok{0.5}\NormalTok{, label}\OperatorTok{=}\StringTok{\textquotesingle{}β₁ verdadero\textquotesingle{}}\NormalTok{)}
\NormalTok{axes[}\DecValTok{0}\NormalTok{].axhline(y}\OperatorTok{=}\NormalTok{beta\_true[}\DecValTok{1}\NormalTok{], color}\OperatorTok{=}\StringTok{\textquotesingle{}orange\textquotesingle{}}\NormalTok{, linestyle}\OperatorTok{=}\StringTok{\textquotesingle{}{-}{-}\textquotesingle{}}\NormalTok{, }
\NormalTok{                alpha}\OperatorTok{=}\FloatTok{0.5}\NormalTok{, label}\OperatorTok{=}\StringTok{\textquotesingle{}β₂ verdadero\textquotesingle{}}\NormalTok{)}
\NormalTok{axes[}\DecValTok{0}\NormalTok{].set\_xlabel(}\StringTok{\textquotesingle{}α (escala logarítmica)\textquotesingle{}}\NormalTok{)}
\NormalTok{axes[}\DecValTok{0}\NormalTok{].set\_ylabel(}\StringTok{\textquotesingle{}Valor del coeficiente\textquotesingle{}}\NormalTok{)}
\NormalTok{axes[}\DecValTok{0}\NormalTok{].set\_title(}\StringTok{\textquotesingle{}Coeficientes vs Parámetro de Regularización\textquotesingle{}}\NormalTok{)}
\NormalTok{axes[}\DecValTok{0}\NormalTok{].legend()}
\NormalTok{axes[}\DecValTok{0}\NormalTok{].grid(}\VariableTok{True}\NormalTok{, alpha}\OperatorTok{=}\FloatTok{0.3}\NormalTok{)}

\CommentTok{\# R²}
\NormalTok{axes[}\DecValTok{1}\NormalTok{].semilogx(df\_ridge[}\StringTok{\textquotesingle{}alpha\textquotesingle{}}\NormalTok{], df\_ridge[}\StringTok{\textquotesingle{}r2\_train\textquotesingle{}}\NormalTok{], }
                \StringTok{\textquotesingle{}o{-}\textquotesingle{}}\NormalTok{, linewidth}\OperatorTok{=}\DecValTok{2}\NormalTok{, markersize}\OperatorTok{=}\DecValTok{8}\NormalTok{, label}\OperatorTok{=}\StringTok{\textquotesingle{}Train\textquotesingle{}}\NormalTok{)}
\NormalTok{axes[}\DecValTok{1}\NormalTok{].semilogx(df\_ridge[}\StringTok{\textquotesingle{}alpha\textquotesingle{}}\NormalTok{], df\_ridge[}\StringTok{\textquotesingle{}r2\_test\textquotesingle{}}\NormalTok{], }
                \StringTok{\textquotesingle{}s{-}\textquotesingle{}}\NormalTok{, linewidth}\OperatorTok{=}\DecValTok{2}\NormalTok{, markersize}\OperatorTok{=}\DecValTok{8}\NormalTok{, label}\OperatorTok{=}\StringTok{\textquotesingle{}Test\textquotesingle{}}\NormalTok{)}
\NormalTok{axes[}\DecValTok{1}\NormalTok{].set\_xlabel(}\StringTok{\textquotesingle{}α (escala logarítmica)\textquotesingle{}}\NormalTok{)}
\NormalTok{axes[}\DecValTok{1}\NormalTok{].set\_ylabel(}\StringTok{\textquotesingle{}R²\textquotesingle{}}\NormalTok{)}
\NormalTok{axes[}\DecValTok{1}\NormalTok{].set\_title(}\StringTok{\textquotesingle{}R² vs Parámetro de Regularización\textquotesingle{}}\NormalTok{)}
\NormalTok{axes[}\DecValTok{1}\NormalTok{].legend()}
\NormalTok{axes[}\DecValTok{1}\NormalTok{].grid(}\VariableTok{True}\NormalTok{, alpha}\OperatorTok{=}\FloatTok{0.3}\NormalTok{)}

\NormalTok{plt.tight\_layout()}
\NormalTok{plt.savefig(}\StringTok{\textquotesingle{}ridge\_regularization.png\textquotesingle{}}\NormalTok{, dpi}\OperatorTok{=}\DecValTok{300}\NormalTok{, bbox\_inches}\OperatorTok{=}\StringTok{\textquotesingle{}tight\textquotesingle{}}\NormalTok{)}
\BuiltInTok{print}\NormalTok{(}\StringTok{"}\CharTok{\textbackslash{}n}\StringTok{Gráfico guardado como \textquotesingle{}ridge\_regularization.png\textquotesingle{}"}\NormalTok{)}

\CommentTok{\# Mejor alpha}
\NormalTok{mejor\_idx }\OperatorTok{=}\NormalTok{ df\_ridge[}\StringTok{\textquotesingle{}r2\_test\textquotesingle{}}\NormalTok{].idxmax()}
\NormalTok{mejor\_alpha }\OperatorTok{=}\NormalTok{ df\_ridge.loc[mejor\_idx, }\StringTok{\textquotesingle{}alpha\textquotesingle{}}\NormalTok{]}
\BuiltInTok{print}\NormalTok{(}\SpecialStringTok{f"}\CharTok{\textbackslash{}n}\SpecialStringTok{Mejor α: }\SpecialCharTok{\{}\NormalTok{mejor\_alpha}\SpecialCharTok{\}}\SpecialStringTok{"}\NormalTok{)}
\BuiltInTok{print}\NormalTok{(}\SpecialStringTok{f"R² Test: }\SpecialCharTok{\{}\NormalTok{df\_ridge}\SpecialCharTok{.}\NormalTok{loc[mejor\_idx, }\StringTok{\textquotesingle{}r2\_test\textquotesingle{}}\NormalTok{]}\SpecialCharTok{:.4f\}}\SpecialStringTok{"}\NormalTok{)}
\end{Highlighting}
\end{Shaded}

\section{Regresión Lasso}\label{regresiuxf3n-lasso}

\subsection{Fundamento matemático}\label{fundamento-matemuxe1tico-2}

\begin{Shaded}
\begin{Highlighting}[]
\BuiltInTok{print}\NormalTok{(}\StringTok{"""}
\StringTok{REGRESIÓN LASSO (L1 REGULARIZATION)}

\StringTok{Objetivo:}
\StringTok{    Minimizar RSS + penalización por valor absoluto de coeficientes}
\StringTok{    }
\StringTok{    min \{ RSS(β) + α·Σ|βⱼ| \}}
\StringTok{    }
\StringTok{    = min \{ Σ(yᵢ {-} ŷᵢ)² + α·Σ|βⱼ| \}}

\StringTok{Donde:}
\StringTok{{-} α: Parámetro de regularización (α ≥ 0)}
\StringTok{{-} ||β||₁: Norma L1 (suma de valores absolutos)}

\StringTok{Diferencia clave con Ridge:}
\StringTok{{-} Lasso puede llevar coeficientes exactamente a CERO}
\StringTok{{-} Ridge solo los reduce pero no los elimina}

\StringTok{Efecto de Lasso:}
\StringTok{✓ Selección automática de variables (feature selection)}
\StringTok{✓ Produce modelos dispersos (sparse models)}
\StringTok{✓ Más interpretable que Ridge}
\StringTok{✓ Útil cuando hay muchas variables irrelevantes}

\StringTok{Cuándo usar Lasso:}
\StringTok{{-} Cuando se cree que solo pocas variables son relevantes}
\StringTok{{-} Cuando se necesita interpretabilidad}
\StringTok{{-} Cuando se quiere selección automática de variables}

\StringTok{Aplicaciones económicas:}
\StringTok{{-} Identificar factores clave que afectan un outcome}
\StringTok{{-} Modelos parsimoniosos para reportes}
\StringTok{{-} Análisis exploratorio de muchas variables}
\StringTok{{-} Predicción con datasets con muchas columnas}
\StringTok{"""}\NormalTok{)}
\end{Highlighting}
\end{Shaded}

\subsection{Implementación}\label{implementaciuxf3n-2}

\begin{Shaded}
\begin{Highlighting}[]
\CommentTok{\# Probar diferentes valores de alpha}
\NormalTok{alphas\_lasso }\OperatorTok{=}\NormalTok{ [}\FloatTok{0.001}\NormalTok{, }\FloatTok{0.01}\NormalTok{, }\FloatTok{0.1}\NormalTok{, }\FloatTok{0.5}\NormalTok{, }\DecValTok{1}\NormalTok{, }\DecValTok{2}\NormalTok{, }\DecValTok{5}\NormalTok{]}

\BuiltInTok{print}\NormalTok{(}\StringTok{"}\CharTok{\textbackslash{}n}\StringTok{REGRESIÓN LASSO {-} EFECTO DE α"}\NormalTok{)}
\BuiltInTok{print}\NormalTok{(}\StringTok{"="} \OperatorTok{*} \DecValTok{80}\NormalTok{)}
\BuiltInTok{print}\NormalTok{(}\SpecialStringTok{f"}\SpecialCharTok{\{}\StringTok{\textquotesingle{}α\textquotesingle{}}\SpecialCharTok{:\textless{}10\}}\SpecialStringTok{ }\SpecialCharTok{\{}\StringTok{\textquotesingle{}β₁\textquotesingle{}}\SpecialCharTok{:\textless{}12\}}\SpecialStringTok{ }\SpecialCharTok{\{}\StringTok{\textquotesingle{}β₂\textquotesingle{}}\SpecialCharTok{:\textless{}12\}}\SpecialStringTok{ }\SpecialCharTok{\{}\StringTok{\textquotesingle{}Variables≠0\textquotesingle{}}\SpecialCharTok{:\textless{}14\}}\SpecialStringTok{ }\SpecialCharTok{\{}\StringTok{\textquotesingle{}R² Train\textquotesingle{}}\SpecialCharTok{:\textless{}12\}}\SpecialStringTok{ }\SpecialCharTok{\{}\StringTok{\textquotesingle{}R² Test\textquotesingle{}}\SpecialCharTok{:\textless{}12\}}\SpecialStringTok{"}\NormalTok{)}
\BuiltInTok{print}\NormalTok{(}\StringTok{"{-}"} \OperatorTok{*} \DecValTok{80}\NormalTok{)}

\NormalTok{resultados\_lasso }\OperatorTok{=}\NormalTok{ []}

\ControlFlowTok{for}\NormalTok{ alpha }\KeywordTok{in}\NormalTok{ alphas\_lasso:}
    \CommentTok{\# Entrenar modelo}
\NormalTok{    modelo\_lasso }\OperatorTok{=}\NormalTok{ Lasso(alpha}\OperatorTok{=}\NormalTok{alpha, max\_iter}\OperatorTok{=}\DecValTok{10000}\NormalTok{)}
\NormalTok{    modelo\_lasso.fit(X\_train, y\_train)}
    
    \CommentTok{\# Contar variables no cero}
\NormalTok{    n\_nonzero }\OperatorTok{=}\NormalTok{ np.}\BuiltInTok{sum}\NormalTok{(modelo\_lasso.coef\_ }\OperatorTok{!=} \DecValTok{0}\NormalTok{)}
    
    \CommentTok{\# Evaluar}
\NormalTok{    r2\_train }\OperatorTok{=}\NormalTok{ modelo\_lasso.score(X\_train, y\_train)}
\NormalTok{    r2\_test }\OperatorTok{=}\NormalTok{ modelo\_lasso.score(X\_test, y\_test)}
    
\NormalTok{    resultados\_lasso.append(\{}
        \StringTok{\textquotesingle{}alpha\textquotesingle{}}\NormalTok{: alpha,}
        \StringTok{\textquotesingle{}beta1\textquotesingle{}}\NormalTok{: modelo\_lasso.coef\_[}\DecValTok{0}\NormalTok{],}
        \StringTok{\textquotesingle{}beta2\textquotesingle{}}\NormalTok{: modelo\_lasso.coef\_[}\DecValTok{1}\NormalTok{],}
        \StringTok{\textquotesingle{}n\_nonzero\textquotesingle{}}\NormalTok{: n\_nonzero,}
        \StringTok{\textquotesingle{}r2\_train\textquotesingle{}}\NormalTok{: r2\_train,}
        \StringTok{\textquotesingle{}r2\_test\textquotesingle{}}\NormalTok{: r2\_test}
\NormalTok{    \})}
    
    \BuiltInTok{print}\NormalTok{(}\SpecialStringTok{f"}\SpecialCharTok{\{}\NormalTok{alpha}\SpecialCharTok{:\textless{}10.3f\}}\SpecialStringTok{ }\SpecialCharTok{\{}\NormalTok{modelo\_lasso}\SpecialCharTok{.}\NormalTok{coef\_[}\DecValTok{0}\NormalTok{]}\SpecialCharTok{:\textless{}12.4f\}}\SpecialStringTok{ "}
          \SpecialStringTok{f"}\SpecialCharTok{\{}\NormalTok{modelo\_lasso}\SpecialCharTok{.}\NormalTok{coef\_[}\DecValTok{1}\NormalTok{]}\SpecialCharTok{:\textless{}12.4f\}}\SpecialStringTok{ }\SpecialCharTok{\{}\NormalTok{n\_nonzero}\SpecialCharTok{:\textless{}14d\}}\SpecialStringTok{ "}
          \SpecialStringTok{f"}\SpecialCharTok{\{}\NormalTok{r2\_train}\SpecialCharTok{:\textless{}12.4f\}}\SpecialStringTok{ }\SpecialCharTok{\{}\NormalTok{r2\_test}\SpecialCharTok{:\textless{}12.4f\}}\SpecialStringTok{"}\NormalTok{)}

\NormalTok{df\_lasso }\OperatorTok{=}\NormalTok{ pd.DataFrame(resultados\_lasso)}
\end{Highlighting}
\end{Shaded}

\subsection{Comparación Ridge vs
Lasso}\label{comparaciuxf3n-ridge-vs-lasso}

\begin{Shaded}
\begin{Highlighting}[]
\CommentTok{\# Visualización comparativa}
\NormalTok{fig, axes }\OperatorTok{=}\NormalTok{ plt.subplots(}\DecValTok{2}\NormalTok{, }\DecValTok{2}\NormalTok{, figsize}\OperatorTok{=}\NormalTok{(}\DecValTok{14}\NormalTok{, }\DecValTok{10}\NormalTok{))}

\CommentTok{\# Coeficiente 1}
\NormalTok{axes[}\DecValTok{0}\NormalTok{, }\DecValTok{0}\NormalTok{].semilogx(df\_ridge[}\StringTok{\textquotesingle{}alpha\textquotesingle{}}\NormalTok{], df\_ridge[}\StringTok{\textquotesingle{}beta1\textquotesingle{}}\NormalTok{], }
                    \StringTok{\textquotesingle{}o{-}\textquotesingle{}}\NormalTok{, linewidth}\OperatorTok{=}\DecValTok{2}\NormalTok{, label}\OperatorTok{=}\StringTok{\textquotesingle{}Ridge\textquotesingle{}}\NormalTok{, markersize}\OperatorTok{=}\DecValTok{8}\NormalTok{)}
\NormalTok{axes[}\DecValTok{0}\NormalTok{, }\DecValTok{0}\NormalTok{].semilogx(df\_lasso[}\StringTok{\textquotesingle{}alpha\textquotesingle{}}\NormalTok{], df\_lasso[}\StringTok{\textquotesingle{}beta1\textquotesingle{}}\NormalTok{], }
                    \StringTok{\textquotesingle{}s{-}\textquotesingle{}}\NormalTok{, linewidth}\OperatorTok{=}\DecValTok{2}\NormalTok{, label}\OperatorTok{=}\StringTok{\textquotesingle{}Lasso\textquotesingle{}}\NormalTok{, markersize}\OperatorTok{=}\DecValTok{8}\NormalTok{)}
\NormalTok{axes[}\DecValTok{0}\NormalTok{, }\DecValTok{0}\NormalTok{].axhline(y}\OperatorTok{=}\NormalTok{beta\_true[}\DecValTok{0}\NormalTok{], color}\OperatorTok{=}\StringTok{\textquotesingle{}black\textquotesingle{}}\NormalTok{, linestyle}\OperatorTok{=}\StringTok{\textquotesingle{}{-}{-}\textquotesingle{}}\NormalTok{, }
\NormalTok{                   alpha}\OperatorTok{=}\FloatTok{0.5}\NormalTok{, label}\OperatorTok{=}\StringTok{\textquotesingle{}Verdadero\textquotesingle{}}\NormalTok{)}
\NormalTok{axes[}\DecValTok{0}\NormalTok{, }\DecValTok{0}\NormalTok{].set\_xlabel(}\StringTok{\textquotesingle{}α\textquotesingle{}}\NormalTok{)}
\NormalTok{axes[}\DecValTok{0}\NormalTok{, }\DecValTok{0}\NormalTok{].set\_ylabel(}\StringTok{\textquotesingle{}β₁\textquotesingle{}}\NormalTok{)}
\NormalTok{axes[}\DecValTok{0}\NormalTok{, }\DecValTok{0}\NormalTok{].set\_title(}\StringTok{\textquotesingle{}Coeficiente β₁ vs α\textquotesingle{}}\NormalTok{)}
\NormalTok{axes[}\DecValTok{0}\NormalTok{, }\DecValTok{0}\NormalTok{].legend()}
\NormalTok{axes[}\DecValTok{0}\NormalTok{, }\DecValTok{0}\NormalTok{].grid(}\VariableTok{True}\NormalTok{, alpha}\OperatorTok{=}\FloatTok{0.3}\NormalTok{)}

\CommentTok{\# Coeficiente 2}
\NormalTok{axes[}\DecValTok{0}\NormalTok{, }\DecValTok{1}\NormalTok{].semilogx(df\_ridge[}\StringTok{\textquotesingle{}alpha\textquotesingle{}}\NormalTok{], df\_ridge[}\StringTok{\textquotesingle{}beta2\textquotesingle{}}\NormalTok{], }
                    \StringTok{\textquotesingle{}o{-}\textquotesingle{}}\NormalTok{, linewidth}\OperatorTok{=}\DecValTok{2}\NormalTok{, label}\OperatorTok{=}\StringTok{\textquotesingle{}Ridge\textquotesingle{}}\NormalTok{, markersize}\OperatorTok{=}\DecValTok{8}\NormalTok{)}
\NormalTok{axes[}\DecValTok{0}\NormalTok{, }\DecValTok{1}\NormalTok{].semilogx(df\_lasso[}\StringTok{\textquotesingle{}alpha\textquotesingle{}}\NormalTok{], df\_lasso[}\StringTok{\textquotesingle{}beta2\textquotesingle{}}\NormalTok{], }
                    \StringTok{\textquotesingle{}s{-}\textquotesingle{}}\NormalTok{, linewidth}\OperatorTok{=}\DecValTok{2}\NormalTok{, label}\OperatorTok{=}\StringTok{\textquotesingle{}Lasso\textquotesingle{}}\NormalTok{, markersize}\OperatorTok{=}\DecValTok{8}\NormalTok{)}
\NormalTok{axes[}\DecValTok{0}\NormalTok{, }\DecValTok{1}\NormalTok{].axhline(y}\OperatorTok{=}\NormalTok{beta\_true[}\DecValTok{1}\NormalTok{], color}\OperatorTok{=}\StringTok{\textquotesingle{}black\textquotesingle{}}\NormalTok{, linestyle}\OperatorTok{=}\StringTok{\textquotesingle{}{-}{-}\textquotesingle{}}\NormalTok{, }
\NormalTok{                   alpha}\OperatorTok{=}\FloatTok{0.5}\NormalTok{, label}\OperatorTok{=}\StringTok{\textquotesingle{}Verdadero\textquotesingle{}}\NormalTok{)}
\NormalTok{axes[}\DecValTok{0}\NormalTok{, }\DecValTok{1}\NormalTok{].set\_xlabel(}\StringTok{\textquotesingle{}α\textquotesingle{}}\NormalTok{)}
\NormalTok{axes[}\DecValTok{0}\NormalTok{, }\DecValTok{1}\NormalTok{].set\_ylabel(}\StringTok{\textquotesingle{}β₂\textquotesingle{}}\NormalTok{)}
\NormalTok{axes[}\DecValTok{0}\NormalTok{, }\DecValTok{1}\NormalTok{].set\_title(}\StringTok{\textquotesingle{}Coeficiente β₂ vs α\textquotesingle{}}\NormalTok{)}
\NormalTok{axes[}\DecValTok{0}\NormalTok{, }\DecValTok{1}\NormalTok{].legend()}
\NormalTok{axes[}\DecValTok{0}\NormalTok{, }\DecValTok{1}\NormalTok{].grid(}\VariableTok{True}\NormalTok{, alpha}\OperatorTok{=}\FloatTok{0.3}\NormalTok{)}

\CommentTok{\# R² Train}
\NormalTok{axes[}\DecValTok{1}\NormalTok{, }\DecValTok{0}\NormalTok{].semilogx(df\_ridge[}\StringTok{\textquotesingle{}alpha\textquotesingle{}}\NormalTok{], df\_ridge[}\StringTok{\textquotesingle{}r2\_train\textquotesingle{}}\NormalTok{], }
                    \StringTok{\textquotesingle{}o{-}\textquotesingle{}}\NormalTok{, linewidth}\OperatorTok{=}\DecValTok{2}\NormalTok{, label}\OperatorTok{=}\StringTok{\textquotesingle{}Ridge\textquotesingle{}}\NormalTok{, markersize}\OperatorTok{=}\DecValTok{8}\NormalTok{)}
\NormalTok{axes[}\DecValTok{1}\NormalTok{, }\DecValTok{0}\NormalTok{].semilogx(df\_lasso[}\StringTok{\textquotesingle{}alpha\textquotesingle{}}\NormalTok{], df\_lasso[}\StringTok{\textquotesingle{}r2\_train\textquotesingle{}}\NormalTok{], }
                    \StringTok{\textquotesingle{}s{-}\textquotesingle{}}\NormalTok{, linewidth}\OperatorTok{=}\DecValTok{2}\NormalTok{, label}\OperatorTok{=}\StringTok{\textquotesingle{}Lasso\textquotesingle{}}\NormalTok{, markersize}\OperatorTok{=}\DecValTok{8}\NormalTok{)}
\NormalTok{axes[}\DecValTok{1}\NormalTok{, }\DecValTok{0}\NormalTok{].set\_xlabel(}\StringTok{\textquotesingle{}α\textquotesingle{}}\NormalTok{)}
\NormalTok{axes[}\DecValTok{1}\NormalTok{, }\DecValTok{0}\NormalTok{].set\_ylabel(}\StringTok{\textquotesingle{}R²\textquotesingle{}}\NormalTok{)}
\NormalTok{axes[}\DecValTok{1}\NormalTok{, }\DecValTok{0}\NormalTok{].set\_title(}\StringTok{\textquotesingle{}R² Entrenamiento vs α\textquotesingle{}}\NormalTok{)}
\NormalTok{axes[}\DecValTok{1}\NormalTok{, }\DecValTok{0}\NormalTok{].legend()}
\NormalTok{axes[}\DecValTok{1}\NormalTok{, }\DecValTok{0}\NormalTok{].grid(}\VariableTok{True}\NormalTok{, alpha}\OperatorTok{=}\FloatTok{0.3}\NormalTok{)}

\CommentTok{\# R² Test}
\NormalTok{axes[}\DecValTok{1}\NormalTok{, }\DecValTok{1}\NormalTok{].semilogx(df\_ridge[}\StringTok{\textquotesingle{}alpha\textquotesingle{}}\NormalTok{], df\_ridge[}\StringTok{\textquotesingle{}r2\_test\textquotesingle{}}\NormalTok{], }
                    \StringTok{\textquotesingle{}o{-}\textquotesingle{}}\NormalTok{, linewidth}\OperatorTok{=}\DecValTok{2}\NormalTok{, label}\OperatorTok{=}\StringTok{\textquotesingle{}Ridge\textquotesingle{}}\NormalTok{, markersize}\OperatorTok{=}\DecValTok{8}\NormalTok{)}
\NormalTok{axes[}\DecValTok{1}\NormalTok{, }\DecValTok{1}\NormalTok{].semilogx(df\_lasso[}\StringTok{\textquotesingle{}alpha\textquotesingle{}}\NormalTok{], df\_lasso[}\StringTok{\textquotesingle{}r2\_test\textquotesingle{}}\NormalTok{], }
                    \StringTok{\textquotesingle{}s{-}\textquotesingle{}}\NormalTok{, linewidth}\OperatorTok{=}\DecValTok{2}\NormalTok{, label}\OperatorTok{=}\StringTok{\textquotesingle{}Lasso\textquotesingle{}}\NormalTok{, markersize}\OperatorTok{=}\DecValTok{8}\NormalTok{)}
\NormalTok{axes[}\DecValTok{1}\NormalTok{, }\DecValTok{1}\NormalTok{].set\_xlabel(}\StringTok{\textquotesingle{}α\textquotesingle{}}\NormalTok{)}
\NormalTok{axes[}\DecValTok{1}\NormalTok{, }\DecValTok{1}\NormalTok{].set\_ylabel(}\StringTok{\textquotesingle{}R²\textquotesingle{}}\NormalTok{)}
\NormalTok{axes[}\DecValTok{1}\NormalTok{, }\DecValTok{1}\NormalTok{].set\_title(}\StringTok{\textquotesingle{}R² Test vs α\textquotesingle{}}\NormalTok{)}
\NormalTok{axes[}\DecValTok{1}\NormalTok{, }\DecValTok{1}\NormalTok{].legend()}
\NormalTok{axes[}\DecValTok{1}\NormalTok{, }\DecValTok{1}\NormalTok{].grid(}\VariableTok{True}\NormalTok{, alpha}\OperatorTok{=}\FloatTok{0.3}\NormalTok{)}

\NormalTok{plt.tight\_layout()}
\NormalTok{plt.savefig(}\StringTok{\textquotesingle{}ridge\_vs\_lasso.png\textquotesingle{}}\NormalTok{, dpi}\OperatorTok{=}\DecValTok{300}\NormalTok{, bbox\_inches}\OperatorTok{=}\StringTok{\textquotesingle{}tight\textquotesingle{}}\NormalTok{)}
\BuiltInTok{print}\NormalTok{(}\StringTok{"}\CharTok{\textbackslash{}n}\StringTok{Gráfico guardado como \textquotesingle{}ridge\_vs\_lasso.png\textquotesingle{}"}\NormalTok{)}
\end{Highlighting}
\end{Shaded}

\section{Elastic Net}\label{elastic-net}

\subsection{Fundamento}\label{fundamento}

\begin{Shaded}
\begin{Highlighting}[]
\BuiltInTok{print}\NormalTok{(}\StringTok{"""}
\StringTok{ELASTIC NET}

\StringTok{Combina penalizaciones L1 (Lasso) y L2 (Ridge):}

\StringTok{    min \{ RSS(β) + α·ρ·Σ|βⱼ| + α·(1{-}ρ)/2·Σβⱼ² \}}

\StringTok{Donde:}
\StringTok{{-} α: Fuerza total de regularización}
\StringTok{{-} ρ: Balance entre L1 y L2 (0 ≤ ρ ≤ 1)}
\StringTok{  * ρ = 0: Ridge puro}
\StringTok{  * ρ = 1: Lasso puro}
\StringTok{  * 0 \textless{} ρ \textless{} 1: Combinación}

\StringTok{Ventajas:}
\StringTok{✓ Combina lo mejor de Ridge y Lasso}
\StringTok{✓ Selección de variables (como Lasso)}
\StringTok{✓ Manejo de variables correlacionadas (como Ridge)}
\StringTok{✓ Más estable que Lasso cuando hay correlaciones altas}

\StringTok{Aplicaciones:}
\StringTok{{-} Cuando hay grupos de variables correlacionadas}
\StringTok{{-} Cuando se necesita selección de variables y estabilidad}
\StringTok{"""}\NormalTok{)}
\end{Highlighting}
\end{Shaded}

\subsection{Implementación}\label{implementaciuxf3n-3}

\begin{Shaded}
\begin{Highlighting}[]
\CommentTok{\# Entrenar Elastic Net}
\NormalTok{modelo\_elastic }\OperatorTok{=}\NormalTok{ ElasticNet(alpha}\OperatorTok{=}\FloatTok{0.1}\NormalTok{, l1\_ratio}\OperatorTok{=}\FloatTok{0.5}\NormalTok{, max\_iter}\OperatorTok{=}\DecValTok{10000}\NormalTok{)}
\NormalTok{modelo\_elastic.fit(X\_train, y\_train)}

\BuiltInTok{print}\NormalTok{(}\StringTok{"}\CharTok{\textbackslash{}n}\StringTok{ELASTIC NET"}\NormalTok{)}
\BuiltInTok{print}\NormalTok{(}\StringTok{"="} \OperatorTok{*} \DecValTok{70}\NormalTok{)}
\BuiltInTok{print}\NormalTok{(}\SpecialStringTok{f"Intercepto: }\SpecialCharTok{\{}\NormalTok{modelo\_elastic}\SpecialCharTok{.}\NormalTok{intercept\_}\SpecialCharTok{:.4f\}}\SpecialStringTok{"}\NormalTok{)}
\BuiltInTok{print}\NormalTok{(}\SpecialStringTok{f"Coeficientes:"}\NormalTok{)}
\BuiltInTok{print}\NormalTok{(}\SpecialStringTok{f"  β₁: }\SpecialCharTok{\{}\NormalTok{modelo\_elastic}\SpecialCharTok{.}\NormalTok{coef\_[}\DecValTok{0}\NormalTok{]}\SpecialCharTok{:.4f\}}\SpecialStringTok{"}\NormalTok{)}
\BuiltInTok{print}\NormalTok{(}\SpecialStringTok{f"  β₂: }\SpecialCharTok{\{}\NormalTok{modelo\_elastic}\SpecialCharTok{.}\NormalTok{coef\_[}\DecValTok{1}\NormalTok{]}\SpecialCharTok{:.4f\}}\SpecialStringTok{"}\NormalTok{)}

\NormalTok{r2\_elastic }\OperatorTok{=}\NormalTok{ modelo\_elastic.score(X\_test, y\_test)}
\BuiltInTok{print}\NormalTok{(}\SpecialStringTok{f"}\CharTok{\textbackslash{}n}\SpecialStringTok{R² Test: }\SpecialCharTok{\{}\NormalTok{r2\_elastic}\SpecialCharTok{:.4f\}}\SpecialStringTok{"}\NormalTok{)}
\end{Highlighting}
\end{Shaded}

\section{Regresión polinomial}\label{regresiuxf3n-polinomial}

\subsection{Fundamento}\label{fundamento-1}

\begin{Shaded}
\begin{Highlighting}[]
\BuiltInTok{print}\NormalTok{(}\StringTok{"""}
\StringTok{REGRESIÓN POLINOMIAL}

\StringTok{Extiende regresión lineal creando características polinomiales:}

\StringTok{Grado 1 (lineal): Y = β₀ + β₁X}
\StringTok{Grado 2: Y = β₀ + β₁X + β₂X²}
\StringTok{Grado 3: Y = β₀ + β₁X + β₂X² + β₃X³}

\StringTok{Proceso:}
\StringTok{1. Crear características polinomiales (X, X², X³, etc.)}
\StringTok{2. Aplicar regresión lineal en características transformadas}

\StringTok{Ventajas:}
\StringTok{✓ Captura relaciones no lineales}
\StringTok{✓ Flexible para ajustar curvas complejas}

\StringTok{Desventajas:}
\StringTok{✗ Propenso a overfitting con grados altos}
\StringTok{✗ Inestable en extremos (extrapolación)}
\StringTok{✗ Difícil de interpretar}

\StringTok{Aplicaciones económicas:}
\StringTok{{-} Curvas de costo (U invertida)}
\StringTok{{-} Funciones de producción con rendimientos decrecientes}
\StringTok{{-} Curvas de demanda no lineales}
\StringTok{{-} Ciclos económicos}
\StringTok{"""}\NormalTok{)}
\end{Highlighting}
\end{Shaded}

\subsection{Implementación}\label{implementaciuxf3n-4}

\begin{Shaded}
\begin{Highlighting}[]
\CommentTok{\# Generar datos con relación no lineal}
\NormalTok{np.random.seed(}\DecValTok{42}\NormalTok{)}
\NormalTok{X\_poly\_1d }\OperatorTok{=}\NormalTok{ np.random.uniform(}\OperatorTok{{-}}\DecValTok{3}\NormalTok{, }\DecValTok{3}\NormalTok{, }\DecValTok{100}\NormalTok{)}
\NormalTok{Y\_poly }\OperatorTok{=} \FloatTok{0.5} \OperatorTok{*}\NormalTok{ X\_poly\_1d}\OperatorTok{**}\DecValTok{2} \OperatorTok{{-}} \DecValTok{2} \OperatorTok{*}\NormalTok{ X\_poly\_1d }\OperatorTok{+} \DecValTok{1} \OperatorTok{+}\NormalTok{ np.random.randn(}\DecValTok{100}\NormalTok{) }\OperatorTok{*} \DecValTok{2}

\CommentTok{\# Ordenar para visualización}
\NormalTok{sort\_idx }\OperatorTok{=}\NormalTok{ np.argsort(X\_poly\_1d)}
\NormalTok{X\_poly\_1d }\OperatorTok{=}\NormalTok{ X\_poly\_1d[sort\_idx]}
\NormalTok{Y\_poly }\OperatorTok{=}\NormalTok{ Y\_poly[sort\_idx]}

\NormalTok{X\_poly\_1d }\OperatorTok{=}\NormalTok{ X\_poly\_1d.reshape(}\OperatorTok{{-}}\DecValTok{1}\NormalTok{, }\DecValTok{1}\NormalTok{)}

\CommentTok{\# Dividir datos}
\NormalTok{X\_train\_poly, X\_test\_poly, y\_train\_poly, y\_test\_poly }\OperatorTok{=}\NormalTok{ train\_test\_split(}
\NormalTok{    X\_poly\_1d, Y\_poly, test\_size}\OperatorTok{=}\FloatTok{0.3}\NormalTok{, random\_state}\OperatorTok{=}\DecValTok{42}
\NormalTok{)}

\CommentTok{\# Probar diferentes grados}
\NormalTok{grados }\OperatorTok{=}\NormalTok{ [}\DecValTok{1}\NormalTok{, }\DecValTok{2}\NormalTok{, }\DecValTok{3}\NormalTok{, }\DecValTok{5}\NormalTok{, }\DecValTok{8}\NormalTok{]}
\NormalTok{colores }\OperatorTok{=}\NormalTok{ [}\StringTok{\textquotesingle{}blue\textquotesingle{}}\NormalTok{, }\StringTok{\textquotesingle{}green\textquotesingle{}}\NormalTok{, }\StringTok{\textquotesingle{}red\textquotesingle{}}\NormalTok{, }\StringTok{\textquotesingle{}purple\textquotesingle{}}\NormalTok{, }\StringTok{\textquotesingle{}orange\textquotesingle{}}\NormalTok{]}

\NormalTok{plt.figure(figsize}\OperatorTok{=}\NormalTok{(}\DecValTok{14}\NormalTok{, }\DecValTok{8}\NormalTok{))}

\BuiltInTok{print}\NormalTok{(}\StringTok{"}\CharTok{\textbackslash{}n}\StringTok{REGRESIÓN POLINOMIAL"}\NormalTok{)}
\BuiltInTok{print}\NormalTok{(}\StringTok{"="} \OperatorTok{*} \DecValTok{70}\NormalTok{)}
\BuiltInTok{print}\NormalTok{(}\SpecialStringTok{f"}\SpecialCharTok{\{}\StringTok{\textquotesingle{}Grado\textquotesingle{}}\SpecialCharTok{:\textless{}8\}}\SpecialStringTok{ }\SpecialCharTok{\{}\StringTok{\textquotesingle{}R² Train\textquotesingle{}}\SpecialCharTok{:\textless{}12\}}\SpecialStringTok{ }\SpecialCharTok{\{}\StringTok{\textquotesingle{}R² Test\textquotesingle{}}\SpecialCharTok{:\textless{}12\}}\SpecialStringTok{ }\SpecialCharTok{\{}\StringTok{\textquotesingle{}MSE Train\textquotesingle{}}\SpecialCharTok{:\textless{}12\}}\SpecialStringTok{ }\SpecialCharTok{\{}\StringTok{\textquotesingle{}MSE Test\textquotesingle{}}\SpecialCharTok{:\textless{}12\}}\SpecialStringTok{"}\NormalTok{)}
\BuiltInTok{print}\NormalTok{(}\StringTok{"{-}"} \OperatorTok{*} \DecValTok{70}\NormalTok{)}

\ControlFlowTok{for}\NormalTok{ grado, color }\KeywordTok{in} \BuiltInTok{zip}\NormalTok{(grados, colores):}
    \CommentTok{\# Crear características polinomiales}
\NormalTok{    poly\_features }\OperatorTok{=}\NormalTok{ PolynomialFeatures(degree}\OperatorTok{=}\NormalTok{grado, include\_bias}\OperatorTok{=}\VariableTok{False}\NormalTok{)}
\NormalTok{    X\_train\_poly\_transformed }\OperatorTok{=}\NormalTok{ poly\_features.fit\_transform(X\_train\_poly)}
\NormalTok{    X\_test\_poly\_transformed }\OperatorTok{=}\NormalTok{ poly\_features.transform(X\_test\_poly)}
    
    \CommentTok{\# Entrenar modelo}
\NormalTok{    modelo }\OperatorTok{=}\NormalTok{ LinearRegression()}
\NormalTok{    modelo.fit(X\_train\_poly\_transformed, y\_train\_poly)}
    
    \CommentTok{\# Evaluar}
\NormalTok{    r2\_train }\OperatorTok{=}\NormalTok{ modelo.score(X\_train\_poly\_transformed, y\_train\_poly)}
\NormalTok{    r2\_test }\OperatorTok{=}\NormalTok{ modelo.score(X\_test\_poly\_transformed, y\_test\_poly)}
    
\NormalTok{    y\_pred\_train }\OperatorTok{=}\NormalTok{ modelo.predict(X\_train\_poly\_transformed)}
\NormalTok{    y\_pred\_test }\OperatorTok{=}\NormalTok{ modelo.predict(X\_test\_poly\_transformed)}
    
\NormalTok{    mse\_train }\OperatorTok{=}\NormalTok{ mean\_squared\_error(y\_train\_poly, y\_pred\_train)}
\NormalTok{    mse\_test }\OperatorTok{=}\NormalTok{ mean\_squared\_error(y\_test\_poly, y\_pred\_test)}
    
    \BuiltInTok{print}\NormalTok{(}\SpecialStringTok{f"}\SpecialCharTok{\{}\NormalTok{grado}\SpecialCharTok{:\textless{}8\}}\SpecialStringTok{ }\SpecialCharTok{\{}\NormalTok{r2\_train}\SpecialCharTok{:\textless{}12.4f\}}\SpecialStringTok{ }\SpecialCharTok{\{}\NormalTok{r2\_test}\SpecialCharTok{:\textless{}12.4f\}}\SpecialStringTok{ "}
          \SpecialStringTok{f"}\SpecialCharTok{\{}\NormalTok{mse\_train}\SpecialCharTok{:\textless{}12.4f\}}\SpecialStringTok{ }\SpecialCharTok{\{}\NormalTok{mse\_test}\SpecialCharTok{:\textless{}12.4f\}}\SpecialStringTok{"}\NormalTok{)}
    
    \CommentTok{\# Graficar}
\NormalTok{    X\_plot }\OperatorTok{=}\NormalTok{ np.linspace(X\_poly\_1d.}\BuiltInTok{min}\NormalTok{(), X\_poly\_1d.}\BuiltInTok{max}\NormalTok{(), }\DecValTok{300}\NormalTok{).reshape(}\OperatorTok{{-}}\DecValTok{1}\NormalTok{, }\DecValTok{1}\NormalTok{)}
\NormalTok{    X\_plot\_transformed }\OperatorTok{=}\NormalTok{ poly\_features.transform(X\_plot)}
\NormalTok{    y\_plot }\OperatorTok{=}\NormalTok{ modelo.predict(X\_plot\_transformed)}
    
\NormalTok{    plt.plot(X\_plot, y\_plot, color}\OperatorTok{=}\NormalTok{color, linewidth}\OperatorTok{=}\DecValTok{2}\NormalTok{, }
\NormalTok{             label}\OperatorTok{=}\SpecialStringTok{f\textquotesingle{}Grado }\SpecialCharTok{\{}\NormalTok{grado}\SpecialCharTok{\}}\SpecialStringTok{ (R²=}\SpecialCharTok{\{}\NormalTok{r2\_test}\SpecialCharTok{:.3f\}}\SpecialStringTok{)\textquotesingle{}}\NormalTok{)}

\CommentTok{\# Datos originales}
\NormalTok{plt.scatter(X\_train\_poly, y\_train\_poly, alpha}\OperatorTok{=}\FloatTok{0.6}\NormalTok{, s}\OperatorTok{=}\DecValTok{50}\NormalTok{, }
\NormalTok{           color}\OperatorTok{=}\StringTok{\textquotesingle{}black\textquotesingle{}}\NormalTok{, label}\OperatorTok{=}\StringTok{\textquotesingle{}Datos entrenamiento\textquotesingle{}}\NormalTok{)}
\NormalTok{plt.scatter(X\_test\_poly, y\_test\_poly, alpha}\OperatorTok{=}\FloatTok{0.6}\NormalTok{, s}\OperatorTok{=}\DecValTok{50}\NormalTok{, }
\NormalTok{           marker}\OperatorTok{=}\StringTok{\textquotesingle{}s\textquotesingle{}}\NormalTok{, color}\OperatorTok{=}\StringTok{\textquotesingle{}gray\textquotesingle{}}\NormalTok{, label}\OperatorTok{=}\StringTok{\textquotesingle{}Datos prueba\textquotesingle{}}\NormalTok{)}

\NormalTok{plt.xlabel(}\StringTok{\textquotesingle{}X\textquotesingle{}}\NormalTok{)}
\NormalTok{plt.ylabel(}\StringTok{\textquotesingle{}Y\textquotesingle{}}\NormalTok{)}
\NormalTok{plt.title(}\StringTok{\textquotesingle{}Regresión Polinomial {-} Efecto del Grado\textquotesingle{}}\NormalTok{)}
\NormalTok{plt.legend()}
\NormalTok{plt.grid(}\VariableTok{True}\NormalTok{, alpha}\OperatorTok{=}\FloatTok{0.3}\NormalTok{)}
\NormalTok{plt.tight\_layout()}
\NormalTok{plt.savefig(}\StringTok{\textquotesingle{}polynomial\_regression.png\textquotesingle{}}\NormalTok{, dpi}\OperatorTok{=}\DecValTok{300}\NormalTok{, bbox\_inches}\OperatorTok{=}\StringTok{\textquotesingle{}tight\textquotesingle{}}\NormalTok{)}
\BuiltInTok{print}\NormalTok{(}\StringTok{"}\CharTok{\textbackslash{}n}\StringTok{Gráfico guardado como \textquotesingle{}polynomial\_regression.png\textquotesingle{}"}\NormalTok{)}
\end{Highlighting}
\end{Shaded}

\section{K-Nearest Neighbors
Regressor}\label{k-nearest-neighbors-regressor}

\subsection{Fundamento}\label{fundamento-2}

\begin{Shaded}
\begin{Highlighting}[]
\BuiltInTok{print}\NormalTok{(}\StringTok{"""}
\StringTok{K{-}NEAREST NEIGHBORS REGRESSOR}

\StringTok{Funcionamiento:}
\StringTok{1. Para predecir Y de un nuevo punto X:}
\StringTok{   a. Encuentra los k vecinos más cercanos en datos de entrenamiento}
\StringTok{   b. Promedia sus valores Y}
\StringTok{   c. Ese promedio es la predicción}

\StringTok{Métricas de distancia:}
\StringTok{{-} Euclidiana: √(Σ(xᵢ {-} yᵢ)²)}
\StringTok{{-} Manhattan: Σ|xᵢ {-} yᵢ|}
\StringTok{{-} Minkowski: (Σ|xᵢ {-} yᵢ|ᵖ)\^{}(1/p)}

\StringTok{Ventajas:}
\StringTok{✓ No asume forma funcional}
\StringTok{✓ Captura relaciones complejas}
\StringTok{✓ Fácil de entender}

\StringTok{Desventajas:}
\StringTok{✗ Lento para predicción con datasets grandes}
\StringTok{✗ Sensible a escala de características}
\StringTok{✗ No funciona bien en alta dimensionalidad}
\StringTok{✗ Necesita elegir k apropiado}

\StringTok{Aplicaciones:}
\StringTok{{-} Predicción de precios de viviendas (casas similares)}
\StringTok{{-} Valuación de activos por comparables}
\StringTok{{-} Forecasting basado en patrones históricos similares}
\StringTok{"""}\NormalTok{)}
\end{Highlighting}
\end{Shaded}

\subsection{Implementación}\label{implementaciuxf3n-5}

\begin{Shaded}
\begin{Highlighting}[]
\CommentTok{\# Estandarizar datos (importante para KNN)}
\NormalTok{scaler }\OperatorTok{=}\NormalTok{ StandardScaler()}
\NormalTok{X\_train\_scaled }\OperatorTok{=}\NormalTok{ scaler.fit\_transform(X\_train)}
\NormalTok{X\_test\_scaled }\OperatorTok{=}\NormalTok{ scaler.transform(X\_test)}

\CommentTok{\# Probar diferentes valores de k}
\NormalTok{k\_values }\OperatorTok{=}\NormalTok{ [}\DecValTok{1}\NormalTok{, }\DecValTok{3}\NormalTok{, }\DecValTok{5}\NormalTok{, }\DecValTok{10}\NormalTok{, }\DecValTok{20}\NormalTok{, }\DecValTok{30}\NormalTok{, }\DecValTok{50}\NormalTok{]}

\BuiltInTok{print}\NormalTok{(}\StringTok{"}\CharTok{\textbackslash{}n}\StringTok{K{-}NEAREST NEIGHBORS REGRESSOR"}\NormalTok{)}
\BuiltInTok{print}\NormalTok{(}\StringTok{"="} \OperatorTok{*} \DecValTok{70}\NormalTok{)}
\BuiltInTok{print}\NormalTok{(}\SpecialStringTok{f"}\SpecialCharTok{\{}\StringTok{\textquotesingle{}k\textquotesingle{}}\SpecialCharTok{:\textless{}8\}}\SpecialStringTok{ }\SpecialCharTok{\{}\StringTok{\textquotesingle{}R² Train\textquotesingle{}}\SpecialCharTok{:\textless{}12\}}\SpecialStringTok{ }\SpecialCharTok{\{}\StringTok{\textquotesingle{}R² Test\textquotesingle{}}\SpecialCharTok{:\textless{}12\}}\SpecialStringTok{ }\SpecialCharTok{\{}\StringTok{\textquotesingle{}MSE Test\textquotesingle{}}\SpecialCharTok{:\textless{}12\}}\SpecialStringTok{"}\NormalTok{)}
\BuiltInTok{print}\NormalTok{(}\StringTok{"{-}"} \OperatorTok{*} \DecValTok{70}\NormalTok{)}

\NormalTok{resultados\_knn }\OperatorTok{=}\NormalTok{ []}

\ControlFlowTok{for}\NormalTok{ k }\KeywordTok{in}\NormalTok{ k\_values:}
    \CommentTok{\# Entrenar modelo}
\NormalTok{    modelo\_knn }\OperatorTok{=}\NormalTok{ KNeighborsRegressor(n\_neighbors}\OperatorTok{=}\NormalTok{k)}
\NormalTok{    modelo\_knn.fit(X\_train\_scaled, y\_train)}
    
    \CommentTok{\# Evaluar}
\NormalTok{    r2\_train }\OperatorTok{=}\NormalTok{ modelo\_knn.score(X\_train\_scaled, y\_train)}
\NormalTok{    r2\_test }\OperatorTok{=}\NormalTok{ modelo\_knn.score(X\_test\_scaled, y\_test)}
    
\NormalTok{    y\_pred\_test }\OperatorTok{=}\NormalTok{ modelo\_knn.predict(X\_test\_scaled)}
\NormalTok{    mse\_test }\OperatorTok{=}\NormalTok{ mean\_squared\_error(y\_test, y\_pred\_test)}
    
\NormalTok{    resultados\_knn.append(\{}
        \StringTok{\textquotesingle{}k\textquotesingle{}}\NormalTok{: k,}
        \StringTok{\textquotesingle{}r2\_train\textquotesingle{}}\NormalTok{: r2\_train,}
        \StringTok{\textquotesingle{}r2\_test\textquotesingle{}}\NormalTok{: r2\_test,}
        \StringTok{\textquotesingle{}mse\_test\textquotesingle{}}\NormalTok{: mse\_test}
\NormalTok{    \})}
    
    \BuiltInTok{print}\NormalTok{(}\SpecialStringTok{f"}\SpecialCharTok{\{}\NormalTok{k}\SpecialCharTok{:\textless{}8\}}\SpecialStringTok{ }\SpecialCharTok{\{}\NormalTok{r2\_train}\SpecialCharTok{:\textless{}12.4f\}}\SpecialStringTok{ }\SpecialCharTok{\{}\NormalTok{r2\_test}\SpecialCharTok{:\textless{}12.4f\}}\SpecialStringTok{ }\SpecialCharTok{\{}\NormalTok{mse\_test}\SpecialCharTok{:\textless{}12.4f\}}\SpecialStringTok{"}\NormalTok{)}

\CommentTok{\# Visualizar}
\NormalTok{df\_knn }\OperatorTok{=}\NormalTok{ pd.DataFrame(resultados\_knn)}

\NormalTok{fig, axes }\OperatorTok{=}\NormalTok{ plt.subplots(}\DecValTok{1}\NormalTok{, }\DecValTok{2}\NormalTok{, figsize}\OperatorTok{=}\NormalTok{(}\DecValTok{14}\NormalTok{, }\DecValTok{5}\NormalTok{))}

\NormalTok{axes[}\DecValTok{0}\NormalTok{].plot(df\_knn[}\StringTok{\textquotesingle{}k\textquotesingle{}}\NormalTok{], df\_knn[}\StringTok{\textquotesingle{}r2\_train\textquotesingle{}}\NormalTok{], }
             \StringTok{\textquotesingle{}o{-}\textquotesingle{}}\NormalTok{, linewidth}\OperatorTok{=}\DecValTok{2}\NormalTok{, label}\OperatorTok{=}\StringTok{\textquotesingle{}Train\textquotesingle{}}\NormalTok{, markersize}\OperatorTok{=}\DecValTok{8}\NormalTok{)}
\NormalTok{axes[}\DecValTok{0}\NormalTok{].plot(df\_knn[}\StringTok{\textquotesingle{}k\textquotesingle{}}\NormalTok{], df\_knn[}\StringTok{\textquotesingle{}r2\_test\textquotesingle{}}\NormalTok{], }
             \StringTok{\textquotesingle{}s{-}\textquotesingle{}}\NormalTok{, linewidth}\OperatorTok{=}\DecValTok{2}\NormalTok{, label}\OperatorTok{=}\StringTok{\textquotesingle{}Test\textquotesingle{}}\NormalTok{, markersize}\OperatorTok{=}\DecValTok{8}\NormalTok{)}
\NormalTok{axes[}\DecValTok{0}\NormalTok{].set\_xlabel(}\StringTok{\textquotesingle{}Número de vecinos (k)\textquotesingle{}}\NormalTok{)}
\NormalTok{axes[}\DecValTok{0}\NormalTok{].set\_ylabel(}\StringTok{\textquotesingle{}R²\textquotesingle{}}\NormalTok{)}
\NormalTok{axes[}\DecValTok{0}\NormalTok{].set\_title(}\StringTok{\textquotesingle{}R² vs k en KNN\textquotesingle{}}\NormalTok{)}
\NormalTok{axes[}\DecValTok{0}\NormalTok{].legend()}
\NormalTok{axes[}\DecValTok{0}\NormalTok{].grid(}\VariableTok{True}\NormalTok{, alpha}\OperatorTok{=}\FloatTok{0.3}\NormalTok{)}

\NormalTok{axes[}\DecValTok{1}\NormalTok{].plot(df\_knn[}\StringTok{\textquotesingle{}k\textquotesingle{}}\NormalTok{], df\_knn[}\StringTok{\textquotesingle{}mse\_test\textquotesingle{}}\NormalTok{], }
             \StringTok{\textquotesingle{}o{-}\textquotesingle{}}\NormalTok{, linewidth}\OperatorTok{=}\DecValTok{2}\NormalTok{, color}\OperatorTok{=}\StringTok{\textquotesingle{}red\textquotesingle{}}\NormalTok{, markersize}\OperatorTok{=}\DecValTok{8}\NormalTok{)}
\NormalTok{axes[}\DecValTok{1}\NormalTok{].set\_xlabel(}\StringTok{\textquotesingle{}Número de vecinos (k)\textquotesingle{}}\NormalTok{)}
\NormalTok{axes[}\DecValTok{1}\NormalTok{].set\_ylabel(}\StringTok{\textquotesingle{}MSE Test\textquotesingle{}}\NormalTok{)}
\NormalTok{axes[}\DecValTok{1}\NormalTok{].set\_title(}\StringTok{\textquotesingle{}MSE vs k en KNN\textquotesingle{}}\NormalTok{)}
\NormalTok{axes[}\DecValTok{1}\NormalTok{].grid(}\VariableTok{True}\NormalTok{, alpha}\OperatorTok{=}\FloatTok{0.3}\NormalTok{)}

\NormalTok{plt.tight\_layout()}
\NormalTok{plt.savefig(}\StringTok{\textquotesingle{}knn\_regression.png\textquotesingle{}}\NormalTok{, dpi}\OperatorTok{=}\DecValTok{300}\NormalTok{, bbox\_inches}\OperatorTok{=}\StringTok{\textquotesingle{}tight\textquotesingle{}}\NormalTok{)}
\BuiltInTok{print}\NormalTok{(}\StringTok{"}\CharTok{\textbackslash{}n}\StringTok{Gráfico guardado como \textquotesingle{}knn\_regression.png\textquotesingle{}"}\NormalTok{)}
\end{Highlighting}
\end{Shaded}

\section{Métricas de evaluación}\label{muxe9tricas-de-evaluaciuxf3n}

\subsection{Explicación de
métricas}\label{explicaciuxf3n-de-muxe9tricas}

\begin{Shaded}
\begin{Highlighting}[]
\BuiltInTok{print}\NormalTok{(}\StringTok{"""}
\StringTok{MÉTRICAS DE EVALUACIÓN EN REGRESIÓN}

\StringTok{1. MSE (Mean Squared Error)}
\StringTok{   MSE = (1/n) Σ(yᵢ {-} ŷᵢ)²}
\StringTok{   }
\StringTok{   {-} Promedio de errores al cuadrado}
\StringTok{   {-} Penaliza más los errores grandes}
\StringTok{   {-} Mismas unidades que Y al cuadrado}
\StringTok{   {-} Siempre positivo, menor es mejor}
\StringTok{   {-} Sensible a outliers}

\StringTok{2. RMSE (Root Mean Squared Error)}
\StringTok{   RMSE = √MSE = √[(1/n) Σ(yᵢ {-} ŷᵢ)²]}
\StringTok{   }
\StringTok{   {-} Raíz cuadrada de MSE}
\StringTok{   {-} Mismas unidades que Y}
\StringTok{   {-} Más interpretable que MSE}
\StringTok{   {-} También sensible a outliers}

\StringTok{3. MAE (Mean Absolute Error)}
\StringTok{   MAE = (1/n) Σ|yᵢ {-} ŷᵢ|}
\StringTok{   }
\StringTok{   {-} Promedio de errores absolutos}
\StringTok{   {-} Mismas unidades que Y}
\StringTok{   {-} Menos sensible a outliers que MSE}
\StringTok{   {-} Todos los errores pesan igual}

\StringTok{4. R² (Coeficiente de determinación)}
\StringTok{   R² = 1 {-} (RSS/TSS) = 1 {-} [Σ(yᵢ {-} ŷᵢ)²] / [Σ(yᵢ {-} ȳ)²]}
\StringTok{   }
\StringTok{   {-} Proporción de varianza explicada}
\StringTok{   {-} Rango: 0 a 1 (puede ser negativo en test)}
\StringTok{   {-} R² = 1: Predicción perfecta}
\StringTok{   {-} R² = 0: Tan bueno como predecir la media}
\StringTok{   {-} No tiene unidades}
\StringTok{   {-} Interpretable e intuitivo}

\StringTok{5. Adjusted R² (R² ajustado)}
\StringTok{   R²ₐdⱼ = 1 {-} [(1 {-} R²)(n {-} 1)] / (n {-} p {-} 1)}
\StringTok{   }
\StringTok{   {-} Penaliza por número de variables}
\StringTok{   {-} Útil para comparar modelos con diferentes \# de variables}
\StringTok{   {-} Puede decrecer si se agregan variables irrelevantes}

\StringTok{Cuál usar:}
\StringTok{{-} R²: Fácil interpretación, comparar modelos}
\StringTok{{-} RMSE: Cuando importan unidades originales}
\StringTok{{-} MAE: Cuando outliers no deberían dominar}
\StringTok{{-} Múltiples métricas: Para evaluación completa}
\StringTok{"""}\NormalTok{)}
\end{Highlighting}
\end{Shaded}

\subsection{Implementación de
métricas}\label{implementaciuxf3n-de-muxe9tricas}

\begin{Shaded}
\begin{Highlighting}[]
\KeywordTok{def}\NormalTok{ calcular\_metricas(y\_true, y\_pred, modelo\_nombre}\OperatorTok{=}\StringTok{"Modelo"}\NormalTok{):}
    \CommentTok{"""Calcula todas las métricas de regresión"""}
    
    \CommentTok{\# Calcular métricas}
\NormalTok{    mse }\OperatorTok{=}\NormalTok{ mean\_squared\_error(y\_true, y\_pred)}
\NormalTok{    rmse }\OperatorTok{=}\NormalTok{ np.sqrt(mse)}
\NormalTok{    mae }\OperatorTok{=}\NormalTok{ mean\_absolute\_error(y\_true, y\_pred)}
\NormalTok{    r2 }\OperatorTok{=}\NormalTok{ r2\_score(y\_true, y\_pred)}
    
    \CommentTok{\# Calcular R² ajustado manualmente}
\NormalTok{    n }\OperatorTok{=} \BuiltInTok{len}\NormalTok{(y\_true)}
\NormalTok{    p }\OperatorTok{=} \DecValTok{2}  \CommentTok{\# número de predictores en nuestro caso}
\NormalTok{    r2\_adj }\OperatorTok{=} \DecValTok{1} \OperatorTok{{-}}\NormalTok{ (}\DecValTok{1} \OperatorTok{{-}}\NormalTok{ r2) }\OperatorTok{*}\NormalTok{ (n }\OperatorTok{{-}} \DecValTok{1}\NormalTok{) }\OperatorTok{/}\NormalTok{ (n }\OperatorTok{{-}}\NormalTok{ p }\OperatorTok{{-}} \DecValTok{1}\NormalTok{)}
    
    \CommentTok{\# Imprimir}
    \BuiltInTok{print}\NormalTok{(}\SpecialStringTok{f"}\CharTok{\textbackslash{}n}\SpecialCharTok{\{}\NormalTok{modelo\_nombre}\SpecialCharTok{\}}\SpecialStringTok{"}\NormalTok{)}
    \BuiltInTok{print}\NormalTok{(}\StringTok{"{-}"} \OperatorTok{*} \DecValTok{50}\NormalTok{)}
    \BuiltInTok{print}\NormalTok{(}\SpecialStringTok{f"  MSE:          }\SpecialCharTok{\{}\NormalTok{mse}\SpecialCharTok{:.4f\}}\SpecialStringTok{"}\NormalTok{)}
    \BuiltInTok{print}\NormalTok{(}\SpecialStringTok{f"  RMSE:         }\SpecialCharTok{\{}\NormalTok{rmse}\SpecialCharTok{:.4f\}}\SpecialStringTok{"}\NormalTok{)}
    \BuiltInTok{print}\NormalTok{(}\SpecialStringTok{f"  MAE:          }\SpecialCharTok{\{}\NormalTok{mae}\SpecialCharTok{:.4f\}}\SpecialStringTok{"}\NormalTok{)}
    \BuiltInTok{print}\NormalTok{(}\SpecialStringTok{f"  R²:           }\SpecialCharTok{\{}\NormalTok{r2}\SpecialCharTok{:.4f\}}\SpecialStringTok{"}\NormalTok{)}
    \BuiltInTok{print}\NormalTok{(}\SpecialStringTok{f"  R² Ajustado:  }\SpecialCharTok{\{}\NormalTok{r2\_adj}\SpecialCharTok{:.4f\}}\SpecialStringTok{"}\NormalTok{)}
    
    \ControlFlowTok{return}\NormalTok{ \{}
        \StringTok{\textquotesingle{}modelo\textquotesingle{}}\NormalTok{: modelo\_nombre,}
        \StringTok{\textquotesingle{}mse\textquotesingle{}}\NormalTok{: mse,}
        \StringTok{\textquotesingle{}rmse\textquotesingle{}}\NormalTok{: rmse,}
        \StringTok{\textquotesingle{}mae\textquotesingle{}}\NormalTok{: mae,}
        \StringTok{\textquotesingle{}r2\textquotesingle{}}\NormalTok{: r2,}
        \StringTok{\textquotesingle{}r2\_adj\textquotesingle{}}\NormalTok{: r2\_adj}
\NormalTok{    \}}

\CommentTok{\# Calcular métricas para cada modelo}
\BuiltInTok{print}\NormalTok{(}\StringTok{"}\CharTok{\textbackslash{}n}\StringTok{MÉTRICAS DE EVALUACIÓN {-} CONJUNTO DE PRUEBA"}\NormalTok{)}
\BuiltInTok{print}\NormalTok{(}\StringTok{"="} \OperatorTok{*} \DecValTok{70}\NormalTok{)}

\NormalTok{metricas\_modelos }\OperatorTok{=}\NormalTok{ []}

\CommentTok{\# OLS}
\NormalTok{metricas\_modelos.append(calcular\_metricas(y\_test, y\_pred\_ols\_test, }\StringTok{"OLS"}\NormalTok{))}

\CommentTok{\# Ridge (con mejor alpha)}
\NormalTok{modelo\_ridge\_final }\OperatorTok{=}\NormalTok{ Ridge(alpha}\OperatorTok{=}\NormalTok{mejor\_alpha)}
\NormalTok{modelo\_ridge\_final.fit(X\_train, y\_train)}
\NormalTok{y\_pred\_ridge }\OperatorTok{=}\NormalTok{ modelo\_ridge\_final.predict(X\_test)}
\NormalTok{metricas\_modelos.append(calcular\_metricas(y\_test, y\_pred\_ridge, }\StringTok{"Ridge"}\NormalTok{))}

\CommentTok{\# Lasso}
\NormalTok{mejor\_alpha\_lasso }\OperatorTok{=}\NormalTok{ df\_lasso.loc[df\_lasso[}\StringTok{\textquotesingle{}r2\_test\textquotesingle{}}\NormalTok{].idxmax(), }\StringTok{\textquotesingle{}alpha\textquotesingle{}}\NormalTok{]}
\NormalTok{modelo\_lasso\_final }\OperatorTok{=}\NormalTok{ Lasso(alpha}\OperatorTok{=}\NormalTok{mejor\_alpha\_lasso, max\_iter}\OperatorTok{=}\DecValTok{10000}\NormalTok{)}
\NormalTok{modelo\_lasso\_final.fit(X\_train, y\_train)}
\NormalTok{y\_pred\_lasso }\OperatorTok{=}\NormalTok{ modelo\_lasso\_final.predict(X\_test)}
\NormalTok{metricas\_modelos.append(calcular\_metricas(y\_test, y\_pred\_lasso, }\StringTok{"Lasso"}\NormalTok{))}

\CommentTok{\# KNN}
\NormalTok{mejor\_k }\OperatorTok{=}\NormalTok{ df\_knn.loc[df\_knn[}\StringTok{\textquotesingle{}r2\_test\textquotesingle{}}\NormalTok{].idxmax(), }\StringTok{\textquotesingle{}k\textquotesingle{}}\NormalTok{]}
\NormalTok{modelo\_knn\_final }\OperatorTok{=}\NormalTok{ KNeighborsRegressor(n\_neighbors}\OperatorTok{=}\BuiltInTok{int}\NormalTok{(mejor\_k))}
\NormalTok{modelo\_knn\_final.fit(X\_train\_scaled, y\_train)}
\NormalTok{y\_pred\_knn }\OperatorTok{=}\NormalTok{ modelo\_knn\_final.predict(X\_test\_scaled)}
\NormalTok{metricas\_modelos.append(calcular\_metricas(y\_test, y\_pred\_knn, }\StringTok{"KNN"}\NormalTok{))}

\CommentTok{\# Crear DataFrame comparativo}
\NormalTok{df\_metricas }\OperatorTok{=}\NormalTok{ pd.DataFrame(metricas\_modelos)}

\BuiltInTok{print}\NormalTok{(}\StringTok{"}\CharTok{\textbackslash{}n\textbackslash{}n}\StringTok{TABLA COMPARATIVA"}\NormalTok{)}
\BuiltInTok{print}\NormalTok{(}\StringTok{"="} \OperatorTok{*} \DecValTok{70}\NormalTok{)}
\BuiltInTok{print}\NormalTok{(df\_metricas.to\_string(index}\OperatorTok{=}\VariableTok{False}\NormalTok{))}
\end{Highlighting}
\end{Shaded}

\section{Comparación de modelos}\label{comparaciuxf3n-de-modelos}

\subsection{Tabla y visualización
comparativa}\label{tabla-y-visualizaciuxf3n-comparativa}

\begin{Shaded}
\begin{Highlighting}[]
\CommentTok{\# Visualización comparativa}
\NormalTok{fig, axes }\OperatorTok{=}\NormalTok{ plt.subplots(}\DecValTok{2}\NormalTok{, }\DecValTok{2}\NormalTok{, figsize}\OperatorTok{=}\NormalTok{(}\DecValTok{14}\NormalTok{, }\DecValTok{10}\NormalTok{))}

\CommentTok{\# R²}
\NormalTok{axes[}\DecValTok{0}\NormalTok{, }\DecValTok{0}\NormalTok{].bar(df\_metricas[}\StringTok{\textquotesingle{}modelo\textquotesingle{}}\NormalTok{], df\_metricas[}\StringTok{\textquotesingle{}r2\textquotesingle{}}\NormalTok{], alpha}\OperatorTok{=}\FloatTok{0.7}\NormalTok{)}
\NormalTok{axes[}\DecValTok{0}\NormalTok{, }\DecValTok{0}\NormalTok{].set\_ylabel(}\StringTok{\textquotesingle{}R²\textquotesingle{}}\NormalTok{)}
\NormalTok{axes[}\DecValTok{0}\NormalTok{, }\DecValTok{0}\NormalTok{].set\_title(}\StringTok{\textquotesingle{}R² por Modelo\textquotesingle{}}\NormalTok{)}
\NormalTok{axes[}\DecValTok{0}\NormalTok{, }\DecValTok{0}\NormalTok{].set\_ylim([}\DecValTok{0}\NormalTok{, }\DecValTok{1}\NormalTok{])}
\NormalTok{axes[}\DecValTok{0}\NormalTok{, }\DecValTok{0}\NormalTok{].grid(}\VariableTok{True}\NormalTok{, alpha}\OperatorTok{=}\FloatTok{0.3}\NormalTok{, axis}\OperatorTok{=}\StringTok{\textquotesingle{}y\textquotesingle{}}\NormalTok{)}

\CommentTok{\# RMSE}
\NormalTok{axes[}\DecValTok{0}\NormalTok{, }\DecValTok{1}\NormalTok{].bar(df\_metricas[}\StringTok{\textquotesingle{}modelo\textquotesingle{}}\NormalTok{], df\_metricas[}\StringTok{\textquotesingle{}rmse\textquotesingle{}}\NormalTok{], }
\NormalTok{              alpha}\OperatorTok{=}\FloatTok{0.7}\NormalTok{, color}\OperatorTok{=}\StringTok{\textquotesingle{}orange\textquotesingle{}}\NormalTok{)}
\NormalTok{axes[}\DecValTok{0}\NormalTok{, }\DecValTok{1}\NormalTok{].set\_ylabel(}\StringTok{\textquotesingle{}RMSE\textquotesingle{}}\NormalTok{)}
\NormalTok{axes[}\DecValTok{0}\NormalTok{, }\DecValTok{1}\NormalTok{].set\_title(}\StringTok{\textquotesingle{}RMSE por Modelo (menor es mejor)\textquotesingle{}}\NormalTok{)}
\NormalTok{axes[}\DecValTok{0}\NormalTok{, }\DecValTok{1}\NormalTok{].grid(}\VariableTok{True}\NormalTok{, alpha}\OperatorTok{=}\FloatTok{0.3}\NormalTok{, axis}\OperatorTok{=}\StringTok{\textquotesingle{}y\textquotesingle{}}\NormalTok{)}

\CommentTok{\# MAE}
\NormalTok{axes[}\DecValTok{1}\NormalTok{, }\DecValTok{0}\NormalTok{].bar(df\_metricas[}\StringTok{\textquotesingle{}modelo\textquotesingle{}}\NormalTok{], df\_metricas[}\StringTok{\textquotesingle{}mae\textquotesingle{}}\NormalTok{], }
\NormalTok{              alpha}\OperatorTok{=}\FloatTok{0.7}\NormalTok{, color}\OperatorTok{=}\StringTok{\textquotesingle{}green\textquotesingle{}}\NormalTok{)}
\NormalTok{axes[}\DecValTok{1}\NormalTok{, }\DecValTok{0}\NormalTok{].set\_ylabel(}\StringTok{\textquotesingle{}MAE\textquotesingle{}}\NormalTok{)}
\NormalTok{axes[}\DecValTok{1}\NormalTok{, }\DecValTok{0}\NormalTok{].set\_title(}\StringTok{\textquotesingle{}MAE por Modelo (menor es mejor)\textquotesingle{}}\NormalTok{)}
\NormalTok{axes[}\DecValTok{1}\NormalTok{, }\DecValTok{0}\NormalTok{].grid(}\VariableTok{True}\NormalTok{, alpha}\OperatorTok{=}\FloatTok{0.3}\NormalTok{, axis}\OperatorTok{=}\StringTok{\textquotesingle{}y\textquotesingle{}}\NormalTok{)}

\CommentTok{\# Comparación múltiple}
\NormalTok{x }\OperatorTok{=}\NormalTok{ np.arange(}\BuiltInTok{len}\NormalTok{(df\_metricas))}
\NormalTok{width }\OperatorTok{=} \FloatTok{0.35}

\NormalTok{axes[}\DecValTok{1}\NormalTok{, }\DecValTok{1}\NormalTok{].bar(x }\OperatorTok{{-}}\NormalTok{ width}\OperatorTok{/}\DecValTok{2}\NormalTok{, df\_metricas[}\StringTok{\textquotesingle{}mse\textquotesingle{}}\NormalTok{], width, }
\NormalTok{              label}\OperatorTok{=}\StringTok{\textquotesingle{}MSE\textquotesingle{}}\NormalTok{, alpha}\OperatorTok{=}\FloatTok{0.7}\NormalTok{)}
\NormalTok{axes[}\DecValTok{1}\NormalTok{, }\DecValTok{1}\NormalTok{].bar(x }\OperatorTok{+}\NormalTok{ width}\OperatorTok{/}\DecValTok{2}\NormalTok{, df\_metricas[}\StringTok{\textquotesingle{}mae\textquotesingle{}}\NormalTok{], width, }
\NormalTok{              label}\OperatorTok{=}\StringTok{\textquotesingle{}MAE\textquotesingle{}}\NormalTok{, alpha}\OperatorTok{=}\FloatTok{0.7}\NormalTok{)}
\NormalTok{axes[}\DecValTok{1}\NormalTok{, }\DecValTok{1}\NormalTok{].set\_xticks(x)}
\NormalTok{axes[}\DecValTok{1}\NormalTok{, }\DecValTok{1}\NormalTok{].set\_xticklabels(df\_metricas[}\StringTok{\textquotesingle{}modelo\textquotesingle{}}\NormalTok{])}
\NormalTok{axes[}\DecValTok{1}\NormalTok{, }\DecValTok{1}\NormalTok{].set\_ylabel(}\StringTok{\textquotesingle{}Error\textquotesingle{}}\NormalTok{)}
\NormalTok{axes[}\DecValTok{1}\NormalTok{, }\DecValTok{1}\NormalTok{].set\_title(}\StringTok{\textquotesingle{}Comparación MSE vs MAE\textquotesingle{}}\NormalTok{)}
\NormalTok{axes[}\DecValTok{1}\NormalTok{, }\DecValTok{1}\NormalTok{].legend()}
\NormalTok{axes[}\DecValTok{1}\NormalTok{, }\DecValTok{1}\NormalTok{].grid(}\VariableTok{True}\NormalTok{, alpha}\OperatorTok{=}\FloatTok{0.3}\NormalTok{, axis}\OperatorTok{=}\StringTok{\textquotesingle{}y\textquotesingle{}}\NormalTok{)}

\NormalTok{plt.tight\_layout()}
\NormalTok{plt.savefig(}\StringTok{\textquotesingle{}model\_comparison.png\textquotesingle{}}\NormalTok{, dpi}\OperatorTok{=}\DecValTok{300}\NormalTok{, bbox\_inches}\OperatorTok{=}\StringTok{\textquotesingle{}tight\textquotesingle{}}\NormalTok{)}
\BuiltInTok{print}\NormalTok{(}\StringTok{"}\CharTok{\textbackslash{}n}\StringTok{Comparación guardada como \textquotesingle{}model\_comparison.png\textquotesingle{}"}\NormalTok{)}

\CommentTok{\# Guardar tabla}
\NormalTok{df\_metricas.to\_csv(}\StringTok{\textquotesingle{}metricas\_comparacion.csv\textquotesingle{}}\NormalTok{, index}\OperatorTok{=}\VariableTok{False}\NormalTok{)}
\BuiltInTok{print}\NormalTok{(}\StringTok{"Tabla guardada como \textquotesingle{}metricas\_comparacion.csv\textquotesingle{}"}\NormalTok{)}
\end{Highlighting}
\end{Shaded}

\section{Aplicaciones económicas}\label{aplicaciones-econuxf3micas}

\subsection{Aplicación 1: Predicción de precios de
vivienda}\label{aplicaciuxf3n-1-predicciuxf3n-de-precios-de-vivienda}

\begin{Shaded}
\begin{Highlighting}[]
\BuiltInTok{print}\NormalTok{(}\StringTok{"}\CharTok{\textbackslash{}n\textbackslash{}n\textbackslash{}n}\StringTok{APLICACIÓN: PREDICCIÓN DE PRECIOS DE VIVIENDA"}\NormalTok{)}
\BuiltInTok{print}\NormalTok{(}\StringTok{"="} \OperatorTok{*} \DecValTok{70}\NormalTok{)}

\CommentTok{\# Generar datos sintéticos de viviendas}
\NormalTok{np.random.seed(}\DecValTok{42}\NormalTok{)}
\NormalTok{n\_casas }\OperatorTok{=} \DecValTok{500}

\NormalTok{datos\_vivienda }\OperatorTok{=}\NormalTok{ pd.DataFrame(\{}
    \StringTok{\textquotesingle{}area\_m2\textquotesingle{}}\NormalTok{: np.random.uniform(}\DecValTok{50}\NormalTok{, }\DecValTok{300}\NormalTok{, n\_casas),}
    \StringTok{\textquotesingle{}num\_habitaciones\textquotesingle{}}\NormalTok{: np.random.randint(}\DecValTok{1}\NormalTok{, }\DecValTok{6}\NormalTok{, n\_casas),}
    \StringTok{\textquotesingle{}num\_banos\textquotesingle{}}\NormalTok{: np.random.randint(}\DecValTok{1}\NormalTok{, }\DecValTok{4}\NormalTok{, n\_casas),}
    \StringTok{\textquotesingle{}antiguedad\_anos\textquotesingle{}}\NormalTok{: np.random.uniform(}\DecValTok{0}\NormalTok{, }\DecValTok{50}\NormalTok{, n\_casas),}
    \StringTok{\textquotesingle{}distancia\_centro\_km\textquotesingle{}}\NormalTok{: np.random.uniform(}\FloatTok{0.5}\NormalTok{, }\DecValTok{30}\NormalTok{, n\_casas),}
    \StringTok{\textquotesingle{}area\_jardin\_m2\textquotesingle{}}\NormalTok{: np.random.uniform(}\DecValTok{0}\NormalTok{, }\DecValTok{200}\NormalTok{, n\_casas),}
    \StringTok{\textquotesingle{}tiene\_cochera\textquotesingle{}}\NormalTok{: np.random.choice([}\DecValTok{0}\NormalTok{, }\DecValTok{1}\NormalTok{], n\_casas, p}\OperatorTok{=}\NormalTok{[}\FloatTok{0.3}\NormalTok{, }\FloatTok{0.7}\NormalTok{]),}
    \StringTok{\textquotesingle{}piso\textquotesingle{}}\NormalTok{: np.random.randint(}\DecValTok{1}\NormalTok{, }\DecValTok{15}\NormalTok{, n\_casas)}
\NormalTok{\})}

\CommentTok{\# Generar precio (en miles de soles)}
\NormalTok{precio }\OperatorTok{=}\NormalTok{ (}
\NormalTok{    datos\_vivienda[}\StringTok{\textquotesingle{}area\_m2\textquotesingle{}}\NormalTok{] }\OperatorTok{*} \DecValTok{5} \OperatorTok{+}
\NormalTok{    datos\_vivienda[}\StringTok{\textquotesingle{}num\_habitaciones\textquotesingle{}}\NormalTok{] }\OperatorTok{*} \DecValTok{50} \OperatorTok{+}
\NormalTok{    datos\_vivienda[}\StringTok{\textquotesingle{}num\_banos\textquotesingle{}}\NormalTok{] }\OperatorTok{*} \DecValTok{30} \OperatorTok{+}
    \OperatorTok{{-}}\NormalTok{datos\_vivienda[}\StringTok{\textquotesingle{}antiguedad\_anos\textquotesingle{}}\NormalTok{] }\OperatorTok{*} \DecValTok{2} \OperatorTok{+}
    \OperatorTok{{-}}\NormalTok{datos\_vivienda[}\StringTok{\textquotesingle{}distancia\_centro\_km\textquotesingle{}}\NormalTok{] }\OperatorTok{*} \DecValTok{8} \OperatorTok{+}
\NormalTok{    datos\_vivienda[}\StringTok{\textquotesingle{}area\_jardin\_m2\textquotesingle{}}\NormalTok{] }\OperatorTok{*} \FloatTok{1.5} \OperatorTok{+}
\NormalTok{    datos\_vivienda[}\StringTok{\textquotesingle{}tiene\_cochera\textquotesingle{}}\NormalTok{] }\OperatorTok{*} \DecValTok{100} \OperatorTok{+}
\NormalTok{    np.random.normal(}\DecValTok{0}\NormalTok{, }\DecValTok{50}\NormalTok{, n\_casas)}
\NormalTok{)}

\NormalTok{datos\_vivienda[}\StringTok{\textquotesingle{}precio\_miles\textquotesingle{}}\NormalTok{] }\OperatorTok{=}\NormalTok{ precio}

\BuiltInTok{print}\NormalTok{(}\SpecialStringTok{f"}\CharTok{\textbackslash{}n}\SpecialStringTok{Dataset de viviendas:"}\NormalTok{)}
\BuiltInTok{print}\NormalTok{(}\SpecialStringTok{f"  Total viviendas: }\SpecialCharTok{\{}\NormalTok{n\_casas}\SpecialCharTok{\}}\SpecialStringTok{"}\NormalTok{)}
\BuiltInTok{print}\NormalTok{(}\SpecialStringTok{f"}\CharTok{\textbackslash{}n}\SpecialStringTok{Estadísticas de precio (miles de soles):"}\NormalTok{)}
\BuiltInTok{print}\NormalTok{(datos\_vivienda[}\StringTok{\textquotesingle{}precio\_miles\textquotesingle{}}\NormalTok{].describe())}

\CommentTok{\# Preparar datos}
\NormalTok{X\_vivienda }\OperatorTok{=}\NormalTok{ datos\_vivienda.drop(}\StringTok{\textquotesingle{}precio\_miles\textquotesingle{}}\NormalTok{, axis}\OperatorTok{=}\DecValTok{1}\NormalTok{)}
\NormalTok{y\_vivienda }\OperatorTok{=}\NormalTok{ datos\_vivienda[}\StringTok{\textquotesingle{}precio\_miles\textquotesingle{}}\NormalTok{]}

\NormalTok{X\_train\_v, X\_test\_v, y\_train\_v, y\_test\_v }\OperatorTok{=}\NormalTok{ train\_test\_split(}
\NormalTok{    X\_vivienda, y\_vivienda, test\_size}\OperatorTok{=}\FloatTok{0.2}\NormalTok{, random\_state}\OperatorTok{=}\DecValTok{42}
\NormalTok{)}

\CommentTok{\# Estandarizar}
\NormalTok{scaler\_v }\OperatorTok{=}\NormalTok{ StandardScaler()}
\NormalTok{X\_train\_v\_scaled }\OperatorTok{=}\NormalTok{ scaler\_v.fit\_transform(X\_train\_v)}
\NormalTok{X\_test\_v\_scaled }\OperatorTok{=}\NormalTok{ scaler\_v.transform(X\_test\_v)}

\CommentTok{\# Entrenar múltiples modelos}
\BuiltInTok{print}\NormalTok{(}\StringTok{"}\CharTok{\textbackslash{}n\textbackslash{}n}\StringTok{RESULTADOS DE MODELOS"}\NormalTok{)}
\BuiltInTok{print}\NormalTok{(}\StringTok{"="} \OperatorTok{*} \DecValTok{70}\NormalTok{)}

\NormalTok{modelos\_vivienda }\OperatorTok{=}\NormalTok{ \{}
    \StringTok{\textquotesingle{}OLS\textquotesingle{}}\NormalTok{: LinearRegression(),}
    \StringTok{\textquotesingle{}Ridge\textquotesingle{}}\NormalTok{: Ridge(alpha}\OperatorTok{=}\DecValTok{10}\NormalTok{),}
    \StringTok{\textquotesingle{}Lasso\textquotesingle{}}\NormalTok{: Lasso(alpha}\OperatorTok{=}\DecValTok{1}\NormalTok{, max\_iter}\OperatorTok{=}\DecValTok{10000}\NormalTok{),}
    \StringTok{\textquotesingle{}Random Forest\textquotesingle{}}\NormalTok{: RandomForestRegressor(n\_estimators}\OperatorTok{=}\DecValTok{100}\NormalTok{, random\_state}\OperatorTok{=}\DecValTok{42}\NormalTok{)}
\NormalTok{\}}

\NormalTok{resultados\_vivienda }\OperatorTok{=}\NormalTok{ []}

\ControlFlowTok{for}\NormalTok{ nombre, modelo }\KeywordTok{in}\NormalTok{ modelos\_vivienda.items():}
    \CommentTok{\# Entrenar (algunos modelos necesitan datos escalados)}
    \ControlFlowTok{if}\NormalTok{ nombre }\KeywordTok{in}\NormalTok{ [}\StringTok{\textquotesingle{}OLS\textquotesingle{}}\NormalTok{, }\StringTok{\textquotesingle{}Ridge\textquotesingle{}}\NormalTok{, }\StringTok{\textquotesingle{}Lasso\textquotesingle{}}\NormalTok{]:}
\NormalTok{        modelo.fit(X\_train\_v\_scaled, y\_train\_v)}
\NormalTok{        y\_pred }\OperatorTok{=}\NormalTok{ modelo.predict(X\_test\_v\_scaled)}
    \ControlFlowTok{else}\NormalTok{:}
\NormalTok{        modelo.fit(X\_train\_v, y\_train\_v)}
\NormalTok{        y\_pred }\OperatorTok{=}\NormalTok{ modelo.predict(X\_test\_v)}
    
    \CommentTok{\# Métricas}
\NormalTok{    r2 }\OperatorTok{=}\NormalTok{ r2\_score(y\_test\_v, y\_pred)}
\NormalTok{    rmse }\OperatorTok{=}\NormalTok{ np.sqrt(mean\_squared\_error(y\_test\_v, y\_pred))}
\NormalTok{    mae }\OperatorTok{=}\NormalTok{ mean\_absolute\_error(y\_test\_v, y\_pred)}
    
\NormalTok{    resultados\_vivienda.append(\{}
        \StringTok{\textquotesingle{}Modelo\textquotesingle{}}\NormalTok{: nombre,}
        \StringTok{\textquotesingle{}R²\textquotesingle{}}\NormalTok{: r2,}
        \StringTok{\textquotesingle{}RMSE\textquotesingle{}}\NormalTok{: rmse,}
        \StringTok{\textquotesingle{}MAE\textquotesingle{}}\NormalTok{: mae}
\NormalTok{    \})}
    
    \BuiltInTok{print}\NormalTok{(}\SpecialStringTok{f"}\CharTok{\textbackslash{}n}\SpecialCharTok{\{}\NormalTok{nombre}\SpecialCharTok{\}}\SpecialStringTok{:"}\NormalTok{)}
    \BuiltInTok{print}\NormalTok{(}\SpecialStringTok{f"  R²:   }\SpecialCharTok{\{}\NormalTok{r2}\SpecialCharTok{:.4f\}}\SpecialStringTok{"}\NormalTok{)}
    \BuiltInTok{print}\NormalTok{(}\SpecialStringTok{f"  RMSE: }\SpecialCharTok{\{}\NormalTok{rmse}\SpecialCharTok{:.2f\}}\SpecialStringTok{ miles de soles"}\NormalTok{)}
    \BuiltInTok{print}\NormalTok{(}\SpecialStringTok{f"  MAE:  }\SpecialCharTok{\{}\NormalTok{mae}\SpecialCharTok{:.2f\}}\SpecialStringTok{ miles de soles"}\NormalTok{)}

\CommentTok{\# Importancia de características (Random Forest)}
\BuiltInTok{print}\NormalTok{(}\SpecialStringTok{f"}\CharTok{\textbackslash{}n\textbackslash{}n}\SpecialStringTok{IMPORTANCIA DE CARACTERÍSTICAS (Random Forest)"}\NormalTok{)}
\BuiltInTok{print}\NormalTok{(}\StringTok{"="} \OperatorTok{*} \DecValTok{70}\NormalTok{)}

\NormalTok{importancias\_vivienda }\OperatorTok{=}\NormalTok{ pd.DataFrame(\{}
    \StringTok{\textquotesingle{}caracteristica\textquotesingle{}}\NormalTok{: X\_vivienda.columns,}
    \StringTok{\textquotesingle{}importancia\textquotesingle{}}\NormalTok{: modelos\_vivienda[}\StringTok{\textquotesingle{}Random Forest\textquotesingle{}}\NormalTok{].feature\_importances\_}
\NormalTok{\}).sort\_values(}\StringTok{\textquotesingle{}importancia\textquotesingle{}}\NormalTok{, ascending}\OperatorTok{=}\VariableTok{False}\NormalTok{)}

\ControlFlowTok{for}\NormalTok{ \_, row }\KeywordTok{in}\NormalTok{ importancias\_vivienda.iterrows():}
\NormalTok{    barra }\OperatorTok{=} \StringTok{\textquotesingle{}█\textquotesingle{}} \OperatorTok{*} \BuiltInTok{int}\NormalTok{(row[}\StringTok{\textquotesingle{}importancia\textquotesingle{}}\NormalTok{] }\OperatorTok{*} \DecValTok{100}\NormalTok{)}
    \BuiltInTok{print}\NormalTok{(}\SpecialStringTok{f"  }\SpecialCharTok{\{}\NormalTok{row[}\StringTok{\textquotesingle{}caracteristica\textquotesingle{}}\NormalTok{]}\SpecialCharTok{:\textless{}25\}}\SpecialStringTok{ }\SpecialCharTok{\{}\NormalTok{row[}\StringTok{\textquotesingle{}importancia\textquotesingle{}}\NormalTok{]}\SpecialCharTok{:.4f\}}\SpecialStringTok{ }\SpecialCharTok{\{}\NormalTok{barra}\SpecialCharTok{\}}\SpecialStringTok{"}\NormalTok{)}

\CommentTok{\# Ejemplo de predicción}
\BuiltInTok{print}\NormalTok{(}\SpecialStringTok{f"}\CharTok{\textbackslash{}n\textbackslash{}n}\SpecialStringTok{EJEMPLO DE PREDICCIÓN"}\NormalTok{)}
\BuiltInTok{print}\NormalTok{(}\StringTok{"="} \OperatorTok{*} \DecValTok{70}\NormalTok{)}

\NormalTok{casa\_ejemplo }\OperatorTok{=}\NormalTok{ pd.DataFrame(\{}
    \StringTok{\textquotesingle{}area\_m2\textquotesingle{}}\NormalTok{: [}\DecValTok{150}\NormalTok{],}
    \StringTok{\textquotesingle{}num\_habitaciones\textquotesingle{}}\NormalTok{: [}\DecValTok{3}\NormalTok{],}
    \StringTok{\textquotesingle{}num\_banos\textquotesingle{}}\NormalTok{: [}\DecValTok{2}\NormalTok{],}
    \StringTok{\textquotesingle{}antiguedad\_anos\textquotesingle{}}\NormalTok{: [}\DecValTok{10}\NormalTok{],}
    \StringTok{\textquotesingle{}distancia\_centro\_km\textquotesingle{}}\NormalTok{: [}\DecValTok{5}\NormalTok{],}
    \StringTok{\textquotesingle{}area\_jardin\_m2\textquotesingle{}}\NormalTok{: [}\DecValTok{50}\NormalTok{],}
    \StringTok{\textquotesingle{}tiene\_cochera\textquotesingle{}}\NormalTok{: [}\DecValTok{1}\NormalTok{],}
    \StringTok{\textquotesingle{}piso\textquotesingle{}}\NormalTok{: [}\DecValTok{3}\NormalTok{]}
\NormalTok{\})}

\BuiltInTok{print}\NormalTok{(}\StringTok{"}\CharTok{\textbackslash{}n}\StringTok{Características de la vivienda:"}\NormalTok{)}
\ControlFlowTok{for}\NormalTok{ col }\KeywordTok{in}\NormalTok{ casa\_ejemplo.columns:}
    \BuiltInTok{print}\NormalTok{(}\SpecialStringTok{f"  }\SpecialCharTok{\{}\NormalTok{col}\SpecialCharTok{\}}\SpecialStringTok{: }\SpecialCharTok{\{}\NormalTok{casa\_ejemplo[col]}\SpecialCharTok{.}\NormalTok{values[}\DecValTok{0}\NormalTok{]}\SpecialCharTok{\}}\SpecialStringTok{"}\NormalTok{)}

\CommentTok{\# Predecir con cada modelo}
\BuiltInTok{print}\NormalTok{(}\SpecialStringTok{f"}\CharTok{\textbackslash{}n}\SpecialStringTok{Predicciones de precio (miles de soles):"}\NormalTok{)}
\NormalTok{casa\_ejemplo\_scaled }\OperatorTok{=}\NormalTok{ scaler\_v.transform(casa\_ejemplo)}

\ControlFlowTok{for}\NormalTok{ nombre, modelo }\KeywordTok{in}\NormalTok{ modelos\_vivienda.items():}
    \ControlFlowTok{if}\NormalTok{ nombre }\KeywordTok{in}\NormalTok{ [}\StringTok{\textquotesingle{}OLS\textquotesingle{}}\NormalTok{, }\StringTok{\textquotesingle{}Ridge\textquotesingle{}}\NormalTok{, }\StringTok{\textquotesingle{}Lasso\textquotesingle{}}\NormalTok{]:}
\NormalTok{        pred }\OperatorTok{=}\NormalTok{ modelo.predict(casa\_ejemplo\_scaled)[}\DecValTok{0}\NormalTok{]}
    \ControlFlowTok{else}\NormalTok{:}
\NormalTok{        pred }\OperatorTok{=}\NormalTok{ modelo.predict(casa\_ejemplo)[}\DecValTok{0}\NormalTok{]}
    
    \BuiltInTok{print}\NormalTok{(}\SpecialStringTok{f"  }\SpecialCharTok{\{}\NormalTok{nombre}\SpecialCharTok{:\textless{}15\}}\SpecialStringTok{ S/ }\SpecialCharTok{\{}\NormalTok{pred}\SpecialCharTok{:,.0f\}}\SpecialStringTok{ mil"}\NormalTok{)}
\end{Highlighting}
\end{Shaded}

\subsection{Aplicación 2: Predicción de
ventas}\label{aplicaciuxf3n-2-predicciuxf3n-de-ventas}

\begin{Shaded}
\begin{Highlighting}[]
\BuiltInTok{print}\NormalTok{(}\StringTok{"}\CharTok{\textbackslash{}n\textbackslash{}n\textbackslash{}n}\StringTok{APLICACIÓN: PREDICCIÓN DE VENTAS"}\NormalTok{)}
\BuiltInTok{print}\NormalTok{(}\StringTok{"="} \OperatorTok{*} \DecValTok{70}\NormalTok{)}

\CommentTok{\# Generar datos de ventas}
\NormalTok{np.random.seed(}\DecValTok{42}\NormalTok{)}
\NormalTok{n\_productos }\OperatorTok{=} \DecValTok{400}

\NormalTok{datos\_ventas }\OperatorTok{=}\NormalTok{ pd.DataFrame(\{}
    \StringTok{\textquotesingle{}precio\textquotesingle{}}\NormalTok{: np.random.uniform(}\DecValTok{10}\NormalTok{, }\DecValTok{200}\NormalTok{, n\_productos),}
    \StringTok{\textquotesingle{}inversion\_publicidad\textquotesingle{}}\NormalTok{: np.random.uniform(}\DecValTok{0}\NormalTok{, }\DecValTok{10000}\NormalTok{, n\_productos),}
    \StringTok{\textquotesingle{}num\_competidores\textquotesingle{}}\NormalTok{: np.random.randint(}\DecValTok{1}\NormalTok{, }\DecValTok{10}\NormalTok{, n\_productos),}
    \StringTok{\textquotesingle{}calificacion\_producto\textquotesingle{}}\NormalTok{: np.random.uniform(}\DecValTok{2}\NormalTok{, }\DecValTok{5}\NormalTok{, n\_productos),}
    \StringTok{\textquotesingle{}inventario\_disponible\textquotesingle{}}\NormalTok{: np.random.randint(}\DecValTok{10}\NormalTok{, }\DecValTok{1000}\NormalTok{, n\_productos),}
    \StringTok{\textquotesingle{}descuento\_pct\textquotesingle{}}\NormalTok{: np.random.uniform(}\DecValTok{0}\NormalTok{, }\DecValTok{30}\NormalTok{, n\_productos),}
    \StringTok{\textquotesingle{}temporada\_alta\textquotesingle{}}\NormalTok{: np.random.choice([}\DecValTok{0}\NormalTok{, }\DecValTok{1}\NormalTok{], n\_productos, p}\OperatorTok{=}\NormalTok{[}\FloatTok{0.7}\NormalTok{, }\FloatTok{0.3}\NormalTok{])}
\NormalTok{\})}

\CommentTok{\# Generar ventas con relación no lineal}
\NormalTok{ventas }\OperatorTok{=}\NormalTok{ (}
    \OperatorTok{{-}}\FloatTok{0.5} \OperatorTok{*}\NormalTok{ datos\_ventas[}\StringTok{\textquotesingle{}precio\textquotesingle{}}\NormalTok{] }\OperatorTok{+}
    \FloatTok{0.8} \OperatorTok{*}\NormalTok{ np.sqrt(datos\_ventas[}\StringTok{\textquotesingle{}inversion\_publicidad\textquotesingle{}}\NormalTok{]) }\OperatorTok{+}
    \OperatorTok{{-}}\DecValTok{10} \OperatorTok{*}\NormalTok{ datos\_ventas[}\StringTok{\textquotesingle{}num\_competidores\textquotesingle{}}\NormalTok{] }\OperatorTok{+}
    \DecValTok{100} \OperatorTok{*}\NormalTok{ datos\_ventas[}\StringTok{\textquotesingle{}calificacion\_producto\textquotesingle{}}\NormalTok{] }\OperatorTok{+}
    \FloatTok{0.05} \OperatorTok{*}\NormalTok{ datos\_ventas[}\StringTok{\textquotesingle{}inventario\_disponible\textquotesingle{}}\NormalTok{] }\OperatorTok{+}
    \DecValTok{3} \OperatorTok{*}\NormalTok{ datos\_ventas[}\StringTok{\textquotesingle{}descuento\_pct\textquotesingle{}}\NormalTok{] }\OperatorTok{+}
    \DecValTok{150} \OperatorTok{*}\NormalTok{ datos\_ventas[}\StringTok{\textquotesingle{}temporada\_alta\textquotesingle{}}\NormalTok{] }\OperatorTok{+}
\NormalTok{    np.random.normal(}\DecValTok{0}\NormalTok{, }\DecValTok{50}\NormalTok{, n\_productos)}
\NormalTok{)}

\NormalTok{datos\_ventas[}\StringTok{\textquotesingle{}ventas\_mensuales\textquotesingle{}}\NormalTok{] }\OperatorTok{=}\NormalTok{ np.maximum(ventas, }\DecValTok{0}\NormalTok{)  }\CommentTok{\# No ventas negativas}

\BuiltInTok{print}\NormalTok{(}\SpecialStringTok{f"}\CharTok{\textbackslash{}n}\SpecialStringTok{Dataset de ventas:"}\NormalTok{)}
\BuiltInTok{print}\NormalTok{(}\SpecialStringTok{f"  Total productos: }\SpecialCharTok{\{}\NormalTok{n\_productos}\SpecialCharTok{\}}\SpecialStringTok{"}\NormalTok{)}
\BuiltInTok{print}\NormalTok{(}\SpecialStringTok{f"}\CharTok{\textbackslash{}n}\SpecialStringTok{Estadísticas de ventas mensuales:"}\NormalTok{)}
\BuiltInTok{print}\NormalTok{(datos\_ventas[}\StringTok{\textquotesingle{}ventas\_mensuales\textquotesingle{}}\NormalTok{].describe())}

\CommentTok{\# Preparar datos}
\NormalTok{X\_ventas }\OperatorTok{=}\NormalTok{ datos\_ventas.drop(}\StringTok{\textquotesingle{}ventas\_mensuales\textquotesingle{}}\NormalTok{, axis}\OperatorTok{=}\DecValTok{1}\NormalTok{)}
\NormalTok{y\_ventas }\OperatorTok{=}\NormalTok{ datos\_ventas[}\StringTok{\textquotesingle{}ventas\_mensuales\textquotesingle{}}\NormalTok{]}

\NormalTok{X\_train\_s, X\_test\_s, y\_train\_s, y\_test\_s }\OperatorTok{=}\NormalTok{ train\_test\_split(}
\NormalTok{    X\_ventas, y\_ventas, test\_size}\OperatorTok{=}\FloatTok{0.2}\NormalTok{, random\_state}\OperatorTok{=}\DecValTok{42}
\NormalTok{)}

\CommentTok{\# Probar regresión polinomial}
\BuiltInTok{print}\NormalTok{(}\StringTok{"}\CharTok{\textbackslash{}n\textbackslash{}n}\StringTok{REGRESIÓN POLINOMIAL PARA VENTAS"}\NormalTok{)}
\BuiltInTok{print}\NormalTok{(}\StringTok{"="} \OperatorTok{*} \DecValTok{70}\NormalTok{)}

\CommentTok{\# Crear características polinomiales de grado 2}
\NormalTok{poly }\OperatorTok{=}\NormalTok{ PolynomialFeatures(degree}\OperatorTok{=}\DecValTok{2}\NormalTok{, include\_bias}\OperatorTok{=}\VariableTok{False}\NormalTok{)}
\NormalTok{X\_train\_s\_poly }\OperatorTok{=}\NormalTok{ poly.fit\_transform(X\_train\_s)}
\NormalTok{X\_test\_s\_poly }\OperatorTok{=}\NormalTok{ poly.transform(X\_test\_s)}

\CommentTok{\# Entrenar modelos}
\NormalTok{modelos\_ventas }\OperatorTok{=}\NormalTok{ \{}
    \StringTok{\textquotesingle{}Lineal\textquotesingle{}}\NormalTok{: LinearRegression(),}
    \StringTok{\textquotesingle{}Polinomial (grado 2)\textquotesingle{}}\NormalTok{: LinearRegression(),}
    \StringTok{\textquotesingle{}Ridge Polinomial\textquotesingle{}}\NormalTok{: Ridge(alpha}\OperatorTok{=}\DecValTok{100}\NormalTok{),}
    \StringTok{\textquotesingle{}Random Forest\textquotesingle{}}\NormalTok{: RandomForestRegressor(n\_estimators}\OperatorTok{=}\DecValTok{200}\NormalTok{, max\_depth}\OperatorTok{=}\DecValTok{10}\NormalTok{, random\_state}\OperatorTok{=}\DecValTok{42}\NormalTok{)}
\NormalTok{\}}

\CommentTok{\# Lineal}
\NormalTok{modelos\_ventas[}\StringTok{\textquotesingle{}Lineal\textquotesingle{}}\NormalTok{].fit(X\_train\_s, y\_train\_s)}
\NormalTok{y\_pred\_lineal }\OperatorTok{=}\NormalTok{ modelos\_ventas[}\StringTok{\textquotesingle{}Lineal\textquotesingle{}}\NormalTok{].predict(X\_test\_s)}

\CommentTok{\# Polinomial}
\NormalTok{modelos\_ventas[}\StringTok{\textquotesingle{}Polinomial (grado 2)\textquotesingle{}}\NormalTok{].fit(X\_train\_s\_poly, y\_train\_s)}
\NormalTok{y\_pred\_poly }\OperatorTok{=}\NormalTok{ modelos\_ventas[}\StringTok{\textquotesingle{}Polinomial (grado 2)\textquotesingle{}}\NormalTok{].predict(X\_test\_s\_poly)}

\CommentTok{\# Ridge Polinomial}
\NormalTok{modelos\_ventas[}\StringTok{\textquotesingle{}Ridge Polinomial\textquotesingle{}}\NormalTok{].fit(X\_train\_s\_poly, y\_train\_s)}
\NormalTok{y\_pred\_ridge\_poly }\OperatorTok{=}\NormalTok{ modelos\_ventas[}\StringTok{\textquotesingle{}Ridge Polinomial\textquotesingle{}}\NormalTok{].predict(X\_test\_s\_poly)}

\CommentTok{\# Random Forest}
\NormalTok{modelos\_ventas[}\StringTok{\textquotesingle{}Random Forest\textquotesingle{}}\NormalTok{].fit(X\_train\_s, y\_train\_s)}
\NormalTok{y\_pred\_rf }\OperatorTok{=}\NormalTok{ modelos\_ventas[}\StringTok{\textquotesingle{}Random Forest\textquotesingle{}}\NormalTok{].predict(X\_test\_s)}

\CommentTok{\# Evaluar}
\NormalTok{predicciones }\OperatorTok{=}\NormalTok{ \{}
    \StringTok{\textquotesingle{}Lineal\textquotesingle{}}\NormalTok{: y\_pred\_lineal,}
    \StringTok{\textquotesingle{}Polinomial (grado 2)\textquotesingle{}}\NormalTok{: y\_pred\_poly,}
    \StringTok{\textquotesingle{}Ridge Polinomial\textquotesingle{}}\NormalTok{: y\_pred\_ridge\_poly,}
    \StringTok{\textquotesingle{}Random Forest\textquotesingle{}}\NormalTok{: y\_pred\_rf}
\NormalTok{\}}

\BuiltInTok{print}\NormalTok{(}\SpecialStringTok{f"}\SpecialCharTok{\{}\StringTok{\textquotesingle{}Modelo\textquotesingle{}}\SpecialCharTok{:\textless{}25\}}\SpecialStringTok{ }\SpecialCharTok{\{}\StringTok{\textquotesingle{}R²\textquotesingle{}}\SpecialCharTok{:\textless{}10\}}\SpecialStringTok{ }\SpecialCharTok{\{}\StringTok{\textquotesingle{}RMSE\textquotesingle{}}\SpecialCharTok{:\textless{}15\}}\SpecialStringTok{ }\SpecialCharTok{\{}\StringTok{\textquotesingle{}MAE\textquotesingle{}}\SpecialCharTok{:\textless{}15\}}\SpecialStringTok{"}\NormalTok{)}
\BuiltInTok{print}\NormalTok{(}\StringTok{"{-}"} \OperatorTok{*} \DecValTok{70}\NormalTok{)}

\ControlFlowTok{for}\NormalTok{ nombre, y\_pred }\KeywordTok{in}\NormalTok{ predicciones.items():}
\NormalTok{    r2 }\OperatorTok{=}\NormalTok{ r2\_score(y\_test\_s, y\_pred)}
\NormalTok{    rmse }\OperatorTok{=}\NormalTok{ np.sqrt(mean\_squared\_error(y\_test\_s, y\_pred))}
\NormalTok{    mae }\OperatorTok{=}\NormalTok{ mean\_absolute\_error(y\_test\_s, y\_pred)}
    
    \BuiltInTok{print}\NormalTok{(}\SpecialStringTok{f"}\SpecialCharTok{\{}\NormalTok{nombre}\SpecialCharTok{:\textless{}25\}}\SpecialStringTok{ }\SpecialCharTok{\{}\NormalTok{r2}\SpecialCharTok{:\textless{}10.4f\}}\SpecialStringTok{ }\SpecialCharTok{\{}\NormalTok{rmse}\SpecialCharTok{:\textless{}15.2f\}}\SpecialStringTok{ }\SpecialCharTok{\{}\NormalTok{mae}\SpecialCharTok{:\textless{}15.2f\}}\SpecialStringTok{"}\NormalTok{)}

\CommentTok{\# Análisis de elasticidad precio}
\BuiltInTok{print}\NormalTok{(}\SpecialStringTok{f"}\CharTok{\textbackslash{}n\textbackslash{}n}\SpecialStringTok{ANÁLISIS DE ELASTICIDAD PRECIO"}\NormalTok{)}
\BuiltInTok{print}\NormalTok{(}\StringTok{"="} \OperatorTok{*} \DecValTok{70}\NormalTok{)}

\CommentTok{\# Usar modelo Random Forest para análisis}
\NormalTok{precios\_test }\OperatorTok{=}\NormalTok{ np.linspace(}\DecValTok{10}\NormalTok{, }\DecValTok{200}\NormalTok{, }\DecValTok{50}\NormalTok{)}
\NormalTok{datos\_base }\OperatorTok{=}\NormalTok{ X\_ventas.mean().to\_dict()}

\NormalTok{ventas\_por\_precio }\OperatorTok{=}\NormalTok{ []}

\ControlFlowTok{for}\NormalTok{ precio }\KeywordTok{in}\NormalTok{ precios\_test:}
\NormalTok{    datos\_pred }\OperatorTok{=}\NormalTok{ pd.DataFrame([datos\_base])}
\NormalTok{    datos\_pred[}\StringTok{\textquotesingle{}precio\textquotesingle{}}\NormalTok{] }\OperatorTok{=}\NormalTok{ precio}
\NormalTok{    ventas\_pred }\OperatorTok{=}\NormalTok{ modelos\_ventas[}\StringTok{\textquotesingle{}Random Forest\textquotesingle{}}\NormalTok{].predict(datos\_pred)[}\DecValTok{0}\NormalTok{]}
\NormalTok{    ventas\_por\_precio.append(ventas\_pred)}

\CommentTok{\# Graficar curva de demanda}
\NormalTok{plt.figure(figsize}\OperatorTok{=}\NormalTok{(}\DecValTok{10}\NormalTok{, }\DecValTok{6}\NormalTok{))}
\NormalTok{plt.plot(precios\_test, ventas\_por\_precio, linewidth}\OperatorTok{=}\DecValTok{2}\NormalTok{)}
\NormalTok{plt.xlabel(}\StringTok{\textquotesingle{}Precio\textquotesingle{}}\NormalTok{)}
\NormalTok{plt.ylabel(}\StringTok{\textquotesingle{}Ventas mensuales predichas\textquotesingle{}}\NormalTok{)}
\NormalTok{plt.title(}\StringTok{\textquotesingle{}Curva de Demanda Estimada\textquotesingle{}}\NormalTok{)}
\NormalTok{plt.grid(}\VariableTok{True}\NormalTok{, alpha}\OperatorTok{=}\FloatTok{0.3}\NormalTok{)}
\NormalTok{plt.tight\_layout()}
\NormalTok{plt.savefig(}\StringTok{\textquotesingle{}curva\_demanda.png\textquotesingle{}}\NormalTok{, dpi}\OperatorTok{=}\DecValTok{300}\NormalTok{, bbox\_inches}\OperatorTok{=}\StringTok{\textquotesingle{}tight\textquotesingle{}}\NormalTok{)}
\BuiltInTok{print}\NormalTok{(}\StringTok{"}\CharTok{\textbackslash{}n}\StringTok{Curva de demanda guardada como \textquotesingle{}curva\_demanda.png\textquotesingle{}"}\NormalTok{)}

\CommentTok{\# Calcular elasticidad en diferentes puntos}
\NormalTok{idx\_medio }\OperatorTok{=} \BuiltInTok{len}\NormalTok{(precios\_test) }\OperatorTok{//} \DecValTok{2}
\NormalTok{precio\_medio }\OperatorTok{=}\NormalTok{ precios\_test[idx\_medio]}
\NormalTok{ventas\_medio }\OperatorTok{=}\NormalTok{ ventas\_por\_precio[idx\_medio]}

\CommentTok{\# Elasticidad = (ΔQ/Q) / (ΔP/P)}
\NormalTok{delta\_p }\OperatorTok{=}\NormalTok{ precios\_test[idx\_medio }\OperatorTok{+} \DecValTok{1}\NormalTok{] }\OperatorTok{{-}}\NormalTok{ precios\_test[idx\_medio]}
\NormalTok{delta\_q }\OperatorTok{=}\NormalTok{ ventas\_por\_precio[idx\_medio }\OperatorTok{+} \DecValTok{1}\NormalTok{] }\OperatorTok{{-}}\NormalTok{ ventas\_por\_precio[idx\_medio]}

\NormalTok{elasticidad }\OperatorTok{=}\NormalTok{ (delta\_q }\OperatorTok{/}\NormalTok{ ventas\_medio) }\OperatorTok{/}\NormalTok{ (delta\_p }\OperatorTok{/}\NormalTok{ precio\_medio)}

\BuiltInTok{print}\NormalTok{(}\SpecialStringTok{f"}\CharTok{\textbackslash{}n}\SpecialStringTok{Elasticidad precio en punto medio:"}\NormalTok{)}
\BuiltInTok{print}\NormalTok{(}\SpecialStringTok{f"  Precio: S/ }\SpecialCharTok{\{}\NormalTok{precio\_medio}\SpecialCharTok{:.2f\}}\SpecialStringTok{"}\NormalTok{)}
\BuiltInTok{print}\NormalTok{(}\SpecialStringTok{f"  Elasticidad: }\SpecialCharTok{\{}\NormalTok{elasticidad}\SpecialCharTok{:.3f\}}\SpecialStringTok{"}\NormalTok{)}

\ControlFlowTok{if} \BuiltInTok{abs}\NormalTok{(elasticidad) }\OperatorTok{\textgreater{}} \DecValTok{1}\NormalTok{:}
    \BuiltInTok{print}\NormalTok{(}\SpecialStringTok{f"  → Demanda elástica (sensible al precio)"}\NormalTok{)}
\ControlFlowTok{else}\NormalTok{:}
    \BuiltInTok{print}\NormalTok{(}\SpecialStringTok{f"  → Demanda inelástica (poco sensible al precio)"}\NormalTok{)}
\end{Highlighting}
\end{Shaded}

\section{Ejercicios prácticos}\label{ejercicios-pruxe1cticos}

\subsection{Ejercicio: Validación
cruzada}\label{ejercicio-validaciuxf3n-cruzada}

\begin{Shaded}
\begin{Highlighting}[]
\BuiltInTok{print}\NormalTok{(}\StringTok{"}\CharTok{\textbackslash{}n\textbackslash{}n\textbackslash{}n}\StringTok{EJERCICIO: VALIDACIÓN CRUZADA"}\NormalTok{)}
\BuiltInTok{print}\NormalTok{(}\StringTok{"="} \OperatorTok{*} \DecValTok{70}\NormalTok{)}

\CommentTok{\# Validación cruzada con diferentes modelos}
\ImportTok{from}\NormalTok{ sklearn.model\_selection }\ImportTok{import}\NormalTok{ cross\_validate}

\NormalTok{modelos\_cv }\OperatorTok{=}\NormalTok{ \{}
    \StringTok{\textquotesingle{}OLS\textquotesingle{}}\NormalTok{: LinearRegression(),}
    \StringTok{\textquotesingle{}Ridge (α=1)\textquotesingle{}}\NormalTok{: Ridge(alpha}\OperatorTok{=}\DecValTok{1}\NormalTok{),}
    \StringTok{\textquotesingle{}Ridge (α=10)\textquotesingle{}}\NormalTok{: Ridge(alpha}\OperatorTok{=}\DecValTok{10}\NormalTok{),}
    \StringTok{\textquotesingle{}Lasso (α=0.1)\textquotesingle{}}\NormalTok{: Lasso(alpha}\OperatorTok{=}\FloatTok{0.1}\NormalTok{, max\_iter}\OperatorTok{=}\DecValTok{10000}\NormalTok{),}
    \StringTok{\textquotesingle{}Lasso (α=1)\textquotesingle{}}\NormalTok{: Lasso(alpha}\OperatorTok{=}\DecValTok{1}\NormalTok{, max\_iter}\OperatorTok{=}\DecValTok{10000}\NormalTok{)}
\NormalTok{\}}

\BuiltInTok{print}\NormalTok{(}\StringTok{"}\CharTok{\textbackslash{}n}\StringTok{Validación cruzada (5{-}fold):"}\NormalTok{)}
\BuiltInTok{print}\NormalTok{(}\SpecialStringTok{f"}\SpecialCharTok{\{}\StringTok{\textquotesingle{}Modelo\textquotesingle{}}\SpecialCharTok{:\textless{}20\}}\SpecialStringTok{ }\SpecialCharTok{\{}\StringTok{\textquotesingle{}R² medio\textquotesingle{}}\SpecialCharTok{:\textless{}12\}}\SpecialStringTok{ }\SpecialCharTok{\{}\StringTok{\textquotesingle{}R² std\textquotesingle{}}\SpecialCharTok{:\textless{}12\}}\SpecialStringTok{ }\SpecialCharTok{\{}\StringTok{\textquotesingle{}RMSE medio\textquotesingle{}}\SpecialCharTok{:\textless{}15\}}\SpecialStringTok{"}\NormalTok{)}
\BuiltInTok{print}\NormalTok{(}\StringTok{"{-}"} \OperatorTok{*} \DecValTok{70}\NormalTok{)}

\NormalTok{resultados\_cv }\OperatorTok{=}\NormalTok{ []}

\ControlFlowTok{for}\NormalTok{ nombre, modelo }\KeywordTok{in}\NormalTok{ modelos\_cv.items():}
    \CommentTok{\# Validación cruzada}
\NormalTok{    cv\_results }\OperatorTok{=}\NormalTok{ cross\_validate(}
\NormalTok{        modelo, X, Y,}
\NormalTok{        cv}\OperatorTok{=}\DecValTok{5}\NormalTok{,}
\NormalTok{        scoring}\OperatorTok{=}\NormalTok{[}\StringTok{\textquotesingle{}r2\textquotesingle{}}\NormalTok{, }\StringTok{\textquotesingle{}neg\_mean\_squared\_error\textquotesingle{}}\NormalTok{],}
\NormalTok{        return\_train\_score}\OperatorTok{=}\VariableTok{False}
\NormalTok{    )}
    
\NormalTok{    r2\_medio }\OperatorTok{=}\NormalTok{ cv\_results[}\StringTok{\textquotesingle{}test\_r2\textquotesingle{}}\NormalTok{].mean()}
\NormalTok{    r2\_std }\OperatorTok{=}\NormalTok{ cv\_results[}\StringTok{\textquotesingle{}test\_r2\textquotesingle{}}\NormalTok{].std()}
\NormalTok{    mse\_medio }\OperatorTok{=} \OperatorTok{{-}}\NormalTok{cv\_results[}\StringTok{\textquotesingle{}test\_neg\_mean\_squared\_error\textquotesingle{}}\NormalTok{].mean()}
\NormalTok{    rmse\_medio }\OperatorTok{=}\NormalTok{ np.sqrt(mse\_medio)}
    
\NormalTok{    resultados\_cv.append(\{}
        \StringTok{\textquotesingle{}modelo\textquotesingle{}}\NormalTok{: nombre,}
        \StringTok{\textquotesingle{}r2\_medio\textquotesingle{}}\NormalTok{: r2\_medio,}
        \StringTok{\textquotesingle{}r2\_std\textquotesingle{}}\NormalTok{: r2\_std,}
        \StringTok{\textquotesingle{}rmse\_medio\textquotesingle{}}\NormalTok{: rmse\_medio}
\NormalTok{    \})}
    
    \BuiltInTok{print}\NormalTok{(}\SpecialStringTok{f"}\SpecialCharTok{\{}\NormalTok{nombre}\SpecialCharTok{:\textless{}20\}}\SpecialStringTok{ }\SpecialCharTok{\{}\NormalTok{r2\_medio}\SpecialCharTok{:\textless{}12.4f\}}\SpecialStringTok{ }\SpecialCharTok{\{}\NormalTok{r2\_std}\SpecialCharTok{:\textless{}12.4f\}}\SpecialStringTok{ }\SpecialCharTok{\{}\NormalTok{rmse\_medio}\SpecialCharTok{:\textless{}15.4f\}}\SpecialStringTok{"}\NormalTok{)}

\CommentTok{\# Mejor modelo por CV}
\NormalTok{df\_cv }\OperatorTok{=}\NormalTok{ pd.DataFrame(resultados\_cv)}
\NormalTok{mejor\_modelo\_cv }\OperatorTok{=}\NormalTok{ df\_cv.loc[df\_cv[}\StringTok{\textquotesingle{}r2\_medio\textquotesingle{}}\NormalTok{].idxmax(), }\StringTok{\textquotesingle{}modelo\textquotesingle{}}\NormalTok{]}

\BuiltInTok{print}\NormalTok{(}\SpecialStringTok{f"}\CharTok{\textbackslash{}n\textbackslash{}n}\SpecialStringTok{Mejor modelo según validación cruzada: }\SpecialCharTok{\{}\NormalTok{mejor\_modelo\_cv}\SpecialCharTok{\}}\SpecialStringTok{"}\NormalTok{)}
\BuiltInTok{print}\NormalTok{(}\StringTok{"(Balance entre performance y generalización)"}\NormalTok{)}
\end{Highlighting}
\end{Shaded}

\section{Conclusión}\label{conclusiuxf3n}

En esta guía hemos explorado los métodos de regresión:

\textbf{Mínimos Cuadrados Ordinarios (OLS)}

Método base, interpretable, con solución analítica. Sensible a outliers
y multicolinealidad.

\textbf{Regresión Ridge}

Penalización L2. Reduce magnitud de coeficientes. Útil con variables
correlacionadas.

\textbf{Regresión Lasso}

Penalización L1. Selección automática de variables. Produce modelos
sparse.

\textbf{Elastic Net}

Combina Ridge y Lasso. Balance entre selección y estabilidad.

\textbf{Regresión Polinomial}

Captura relaciones no lineales. Propenso a overfitting con grados altos.

\textbf{K-Nearest Neighbors}

No paramétrico. Captura patrones complejos. Sensible a escala y lento.

\textbf{Métricas: R², MSE, RMSE, MAE}

Diferentes perspectivas del error. R² para varianza explicada. RMSE/MAE
para magnitud del error.

\section{Próximos pasos}\label{pruxf3ximos-pasos}

En la siguiente guía (Guía 11: Validación Cruzada y Composición del
Modelo) exploraremos:

\begin{itemize}
\tightlist
\item
  Validación cruzada (k-fold, stratified)
\item
  Grid search para hiperparámetros
\item
  Curvas de aprendizaje
\item
  Contribución de features
\item
  Bias-variance tradeoff
\end{itemize}

\section{Recursos adicionales}\label{recursos-adicionales}

Para profundizar en regresión:

\begin{itemize}
\tightlist
\item
  Scikit-learn Regression:
  scikit-learn.org/stable/supervised\_learning.html\#supervised-learning
\item
  Introduction to Statistical Learning (ISLR) - Capítulos 3, 6
\item
  Elements of Statistical Learning (ESL)
\item
  Applied Econometrics with Python
\end{itemize}

\section{Publicaciones Similares}\label{publicaciones-similares}

Si te interesó este artículo, te recomendamos que explores otros blogs y
recursos relacionados que pueden ampliar tus conocimientos. Aquí te dejo
algunas sugerencias:

\begin{enumerate}
\def\labelenumi{\arabic{enumi}.}
\tightlist
\item
  \href{https://numerus-scriptum.netlify.app/python/2020-06-19-instalacion-de-anaconda/index.pdf}{\faIcon{file-pdf}}
  \href{https://numerus-scriptum.netlify.app/python/2020-06-19-instalacion-de-anaconda}{Instalacion
  De Anaconda}
\item
  \href{https://numerus-scriptum.netlify.app/python/2020-06-20-configurar-entorno-virtual-python-anaconda/index.pdf}{\faIcon{file-pdf}}
  \href{https://numerus-scriptum.netlify.app/python/2020-06-20-configurar-entorno-virtual-python-anaconda}{Configurar
  Entorno Virtual Python Anaconda}
\item
  \href{https://numerus-scriptum.netlify.app/python/2021-04-17-01-introducion-a-la-programacion-con-python/index.pdf}{\faIcon{file-pdf}}
  \href{https://numerus-scriptum.netlify.app/python/2021-04-17-01-introducion-a-la-programacion-con-python}{01
  Introducion A La Programacion Con Python}
\item
  \href{https://numerus-scriptum.netlify.app/python/2021-05-31-02-variables-expresiones-y-statements-con-python/index.pdf}{\faIcon{file-pdf}}
  \href{https://numerus-scriptum.netlify.app/python/2021-05-31-02-variables-expresiones-y-statements-con-python}{02
  Variables Expresiones Y Statements Con Python}
\item
  \href{https://numerus-scriptum.netlify.app/python/2021-06-07-03-objetos-de-python/index.pdf}{\faIcon{file-pdf}}
  \href{https://numerus-scriptum.netlify.app/python/2021-06-07-03-objetos-de-python}{03
  Objetos De Python}
\item
  \href{https://numerus-scriptum.netlify.app/python/2021-06-14-04-ejecucion-condicional-con-python/index.pdf}{\faIcon{file-pdf}}
  \href{https://numerus-scriptum.netlify.app/python/2021-06-14-04-ejecucion-condicional-con-python}{04
  Ejecucion Condicional Con Python}
\item
  \href{https://numerus-scriptum.netlify.app/python/2021-06-21-05-iteraciones-con-python/index.pdf}{\faIcon{file-pdf}}
  \href{https://numerus-scriptum.netlify.app/python/2021-06-21-05-iteraciones-con-python}{05
  Iteraciones Con Python}
\item
  \href{https://numerus-scriptum.netlify.app/python/2021-08-16-06-funciones-con-python/index.pdf}{\faIcon{file-pdf}}
  \href{https://numerus-scriptum.netlify.app/python/2021-08-16-06-funciones-con-python}{06
  Funciones Con Python}
\item
  \href{https://numerus-scriptum.netlify.app/python/2021-08-23-07-dataframes-con-python/index.pdf}{\faIcon{file-pdf}}
  \href{https://numerus-scriptum.netlify.app/python/2021-08-23-07-dataframes-con-python}{07
  Dataframes Con Python}
\item
  \href{https://numerus-scriptum.netlify.app/python/2021-11-29-08-prediccion-y-metrica-de-performance-con-python/index.pdf}{\faIcon{file-pdf}}
  \href{https://numerus-scriptum.netlify.app/python/2021-11-29-08-prediccion-y-metrica-de-performance-con-python}{08
  Prediccion Y Metrica De Performance Con Python}
\item
  \href{https://numerus-scriptum.netlify.app/python/2021-12-06-09-metodos-de-machine-learning-para-clasificacion-con-python/index.pdf}{\faIcon{file-pdf}}
  \href{https://numerus-scriptum.netlify.app/python/2021-12-06-09-metodos-de-machine-learning-para-clasificacion-con-python}{09
  Metodos De Machine Learning Para Clasificacion Con Python}
\item
  \href{https://numerus-scriptum.netlify.app/python/2021-12-13-10-metodos-de-machine-learning-para-regresion-con-python/index.pdf}{\faIcon{file-pdf}}
  \href{https://numerus-scriptum.netlify.app/python/2021-12-13-10-metodos-de-machine-learning-para-regresion-con-python}{10
  Metodos De Machine Learning Para Regresion Con Python}
\item
  \href{https://numerus-scriptum.netlify.app/python/2022-10-31-11-validacion-cruzada-y-composicion-del-modelo-con-python/index.pdf}{\faIcon{file-pdf}}
  \href{https://numerus-scriptum.netlify.app/python/2022-10-31-11-validacion-cruzada-y-composicion-del-modelo-con-python}{11
  Validacion Cruzada Y Composicion Del Modelo Con Python}
\item
  \href{https://numerus-scriptum.netlify.app/python/2025-05-10-visualizacion-de-datos-con-python/index.pdf}{\faIcon{file-pdf}}
  \href{https://numerus-scriptum.netlify.app/python/2025-05-10-visualizacion-de-datos-con-python}{Visualizacion
  De Datos Con Python}
\end{enumerate}

Esperamos que encuentres estas publicaciones igualmente interesantes y
útiles. ¡Disfruta de la lectura!






\end{document}
